\section{Definitions and properties}
\label{sec:loosening-defs}

We begin by pinning down \emph{which} types may be loosened into which, before turning, in \Cref{sec:loosening-rules}, to the construction that actually performs the loosening. Two notions serve this purpose: a relation $\sqsubseteq$ stating that a type \emph{can} be loosened into another, and an inductive datatype $\castpath{}{}$ whose inhabitants are explicit witnesses of \emph{how}. The former is convenient for stating properties; the latter, being a concrete piece of data, is what the recursive construction consumes. We first define both, then relate them and study the order-theoretic structure $\sqsubseteq$ induces on $\mathsf{SMTType}$---a structure rich enough to admit normal forms and a Galois connection~\cite{davey-priestley}, which together explain the geometry of loosening.

\subsection{Definitions}

The two definitions share their inference rules: $\regnotation{\castable{\alpha}{\beta}}$ is the proposition ``$\alpha$ may be loosened into $\beta$'', or equivalently, ``$\beta$ is looser than $\alpha$'', whereas $\regnotation{\castpath{\alpha}{\beta}}$ is the type of derivations of that fact. Each rule corresponds to one way of relaxing a representation.

\begin{definition}[Castable relation][castable?][Encoder/Loosening/Castable.html\#castable?]
  \label{def:castable}
  The \emph{castable} relation $\sqsubseteq$ is defined inductively, as follows:
\begin{mathpar}
\typingrule{castable}{refl}
  {\alpha = \intS \lor \alpha = \boolS \lor \alpha = \unitS}
  {\castable{\alpha}{\alpha}}
\and
\typingrule{castable}{opt}
  {\castable{\alpha}{\alpha'}}
  {\castable{(\optS\ \alpha)}{(\optS\ \alpha')}}
\and
\typingrule{castable}{chpred}
  {\castable{\alpha}{\alpha'}}
  {\castable{(\alpha \toS \boolS)}{(\alpha' \toS \boolS)}}
\and
\typingrule{castable}{pair}
  {\castable{\alpha}{\alpha'} \quad \castable{\beta}{\beta'}}
  {\castable{(\alpha \timesS \beta)}{(\alpha' \timesS \beta')}}
\and
\typingrule{castable}{graph}
  {\castable{\alpha}{\alpha'} \quad \castable{\beta}{\beta'}}
  {\castable{(\alpha \toS \optS\ \beta)}{(\alpha' \timesS \beta') \toS \boolS}}
\and
\typingrule{castable}{fun}
  {\castable{\alpha}{\alpha'} \quad \castable{\beta}{\beta'} \quad \beta \neq \boolS}
  {\castable{(\alpha \toS \beta)}{(\alpha' \toS \beta')}}
\end{mathpar}
\end{definition}

The base case \ref{rule:castable-refl} keeps atomic types unchanged. The structural rules \ref{rule:castable-opt}, \ref{rule:castable-chpred}, \ref{rule:castable-pair} and \ref{rule:castable-fun} loosen componentwise under an option, a characteristic-predicate domain, a product, and a function space. The essential rule \ref{rule:castable-graph} turns the functional encoding $\alpha \toS \optS\,\beta$ of a partial function into the relational encoding $(\alpha' \timesS \beta') \toS \boolS$ of its graph---the move motivated in \Cref{sec:loosening-motivation}. The side condition $\beta \neq \boolS$ makes \ref{rule:castable-fun} and \ref{rule:castable-chpred} disjoint: a function into $\boolS$ is a characteristic predicate and is handled only by \ref{rule:castable-chpred}.

The cast path $\castpath{}{}$ replays the very same rules, but as a \emph{proof-relevant}, or \emph{constructive} version of them: where $\castable{\alpha}{\beta}$ is a mere proposition (it lives in \leaninl{Prop}), a cast path $\castpath{\alpha}{\beta}$ is a piece of data (it lives in \leaninl{Type}) recording \emph{how} $\alpha$ is loosened into $\beta$.

\begin{definition}[Cast path][castPath][Encoder/Loosening/Castable.html\#castPath]
  \label{def:castpath}
  The cast path $\castpath{}{}$ is defined inductively, as follows:
\begin{mathpar}
\typingrule{castpath}{refl}
  {\alpha = \intS \lor \alpha = \boolS \lor \alpha = \unitS}
  {\castpath{\alpha}{\alpha}}
\and
\typingrule{castpath}{opt}
  {\castpath{\alpha}{\alpha'}}
  {\castpath{(\optS\ \alpha)}{(\optS\ \alpha')}}
\and
\typingrule{castpath}{chpred}
  {\castpath{\alpha}{\alpha'}}
  {\castpath{(\alpha \toS \boolS)}{(\alpha' \toS \boolS)}}
\and
\typingrule{castpath}{pair}
  {\castpath{\alpha}{\alpha'} \quad \castpath{\beta}{\beta'}}
  {\castpath{(\alpha \timesS \beta)}{(\alpha' \timesS \beta')}}
\and
\typingrule{castpath}{graph}
  {\castpath{\alpha}{\alpha'} \quad \castpath{\beta}{\beta'}}
  {\castpath{(\alpha \toS \optS\ \beta)}{(\alpha' \timesS \beta') \toS \boolS}}
\and
\typingrule{castpath}{fun}
  {\castpath{\alpha}{\alpha'} \quad \castpath{\beta}{\beta'} \quad \beta \neq \boolS}
  {\castpath{(\alpha \toS \beta)}{(\alpha' \toS \beta')}}
\end{mathpar}
\end{definition}
For the reason detailed in \Cref{par:large-elimination}, the two cannot be factored out into a single universe-polymorphic inductive type \leaninl|castable.{u} : SMTType → SMTType → Sort u|. We therefore keep both, and the two agree, as recorded by \Cref{cor:castable-constructive} below.

\begin{remark}
One might be tempted to merge \ref{rule:castable-chpred} into \ref{rule:castable-fun} by dropping the side condition $\beta \neq \boolS$: by inversion, $\castable{\boolS}{\gamma}$ forces $\gamma = \boolS$, so the merged rule would derive exactly the same relation $\sqsubseteq$. We nevertheless keep the two rules separate. The construction of \Cref{sec:loosening-rules} recurses on cast-path constructors, and the \texttt{chpred} equation emits the direct, comprehension-style specification appropriate for sets (\Cref{case:chpred}), whereas a merged \texttt{fun} equation would emit the heavier pointwise form with default values (\Cref{case:fun}). Keeping the constructors disjoint also keeps cast paths canonical, a property the mechanization relies on.
\end{remark}

\begin{example}
The functional encoding $\intS \toS \optS\ \intS$ of a partial function on integers is loosened to the relational encoding $(\intS \timesS \intS) \toS \boolS$ of its graph by a single \texttt{graph} step over two reflexive sub-derivations:
\[
\infer*[right=\texttt{graph}]{
  \infer*[right=\texttt{refl}]{ }{\castpath{\intS}{\intS}}
  \quad
  \infer*[right=\texttt{refl}]{ }{\castpath{\intS}{\intS}}
}{\castpath{(\intS \toS \optS\ \intS)}{(\intS \timesS \intS) \toS \boolS}}
\]
\end{example}

\subsection{Basic properties}

Let $\alpha$, $\beta$, and $\delta$ be SMT types. The relation $\castable{}{}$ enjoys the properties one would expect of a partial order, and---crucially for the construction to come---it is interchangeable with its constructive counterpart $\castpath{}{}$.

\begin{theorem}[Decidability][instDecidableCastable?][Encoder/Loosening/Castable.html\#instDecidableCastable?]
  \label{thm:castable-decidable}
  $\castable{\alpha}{\beta}$ is decidable.
\end{theorem}
\begin{proof}
By structural recursion on $\alpha$ and $\beta$: each rule is selected by the head constructors of $\alpha$ and $\beta$, decidable equality on $\mathsf{SMTType}$ settles the atomic \ref{rule:castable-refl} case, and the structural cases reduce to deciding loosenability of the components. The argument is routine but lengthy, as it enumerates every case explicitly.
\end{proof}

Decidability is what makes loosening usable in the encoder, \ie at runtime: confronted with two SMT types, it can test whether one loosens into the other and, if so, recover an explicit witness. The~proof~of~\Cref{thm:castable-decidable} provides the discrimination tree for the decision procedure.
The next two corollaries record this passage between the propositional and the proof-relevant views.

\begin{corollary}[Constructiveness and convertibility][castable?.toCastPath, castable?\_of\_castPath][
  Encoder/Loosening/Castable.html\#castable?.toCastPath,
  Encoder/Loosening/Castable.html\#castable?_of_castPath]
\label{cor:castable-constructive}
  If $\castable{\alpha}{\beta}$, a witness $\mathbb{c} : \castpath{\alpha}{\beta}$ can be constructed.
  Conversely, if $\castpath{\alpha}{\beta}$ is inhabited, then $\castable{\alpha}{\beta}$.
\end{corollary}

\begin{proof}
  The decision procedure of the previous theorem already builds a cast path whenever it succeeds, so a proof of $\castable{\alpha}{\beta}$ yields one by re-running it---this is precisely how the construction sidesteps the impossibility of eliminating the \leaninl{Prop} directly. Conversely, erasing the data of a cast path---keeping only \emph{that} it exists, not \emph{how}---sends each constructor to the homonymous rule of $\castable{}{}$.
\end{proof}

Together they let us conflate the two notions freely: we write $\castable{\alpha}{\beta}$ when only the fact matters, and pass to a cast path $\mathbb{c} : \castpath{\alpha}{\beta}$ whenever the construction must recurse on the shape of the derivation. Finally, loosening is a genuine ordering, so a type has a well-defined position relative to all those it can be loosened into.

\begin{theorem}[Partial ordering][castable?.reflexive, castable?.trans, castable?.antisymm][Encoder/Loosening/Castable.html\#castable?.reflexive, Encoder/Loosening/Castable.html\#castable?.trans, Encoder/Loosening/Castable.html\#castable?.antisymm]
  \label{thm:loosening-partial-order}
  The relation $\castable{}{}$ is a partial order on $\mathsf{SMTType}$.
\end{theorem}
\begin{proof}
Reflexivity holds by a structural recursion applying \cref{rule:castable-refl} at the leaves and the matching structural rule at each connective. Transitivity is an induction on the two derivations, composing them rule by rule. Antisymmetry is an induction on the derivations of $\castable{\alpha}{\beta}$ and $\castable{\beta}{\alpha}$: the structural rules act componentwise, so the inductive hypotheses give equality of the components, while the shape-changing \cref{rule:castable-graph} admits no converse and thus cannot occur in both directions.
\end{proof}

\begin{corollary}[][castPath.reflexive, castPath.trans, castPath.antisymm][Encoder/Loosening/Castable.html\#castPath.reflexive, Encoder/Loosening/Castable.html\#castPath.trans, Encoder/Loosening/Castable.html\#castPath.antisymm]
  A witness \(\mathbb{r} : \castpath{\alpha}{\alpha}\) can be constructed for any type \(\alpha\).
  From $\mathbb{c}_1 : \castpath{\alpha}{\beta}$ and $\mathbb{c}_2 : \castpath{\beta}{\delta}$, a witness $\mathbb{c} : \castpath{\alpha}{\delta}$ can be constructed; and from $\mathbb{c}_1 : \castpath{\alpha}{\beta}$ and $\mathbb{c}_2 : \castpath{\beta}{\alpha}$, we show that $\alpha = \beta$.
\end{corollary}
\begin{proof}
The same inductions, carried out on cast paths rather than on proofs: reflexivity follows by structural recursion,
transitivity concatenates the derivations into a single path, and antisymmetry establishes $\alpha = \beta$.
\end{proof}

In sum, $\castable{}{}$ and $\castpath{}{}$ present one and the same partial order---once as a proposition, once as data. We now turn to its order-theoretic structure, which is richer than the bare ordering suggests.

\subsection{Normal forms}
\label{sec:loosening-normal-forms}

The structure induced by $\castable{}{}$ is neither a lattice nor boundedly complete (also called Dedekind complete), since $\sqsubseteq$ does not have meets or joins in general (\eg $\intS$ and $\boolS$).
However, every SMT type $\alpha$ admits two normal forms called the \emph{least} and \emph{most general forms} of $\alpha$.

\begin{definition}[Most general form][mgf][Encoder/Loosening/Castable.html\#SMT.SMTType.mgf]
  \label{def:mgf}
    The most general form of $\alpha : \mathsf{SMTType}$, denoted $\regnotation{\mgf{\alpha}}$ or $\mathsf{mgf}(\alpha)$, is defined inductively on $\alpha$, as follows:
  \begin{align*}
    \mgf{\alpha} &\triangleq \alpha &&\text{if \(\alpha\) is one of } \intS, \boolS, \unitS\\
    \mgf{\alpha} &\triangleq (\mgf{\gamma} \timesS \mgf{\delta}) \toS \boolS &&\text{if \(\alpha = \gamma \toS \optS\ \delta\)}
  \end{align*} 
  Remaining cases are defined by structural recursion on $\alpha$.
\end{definition}
%
Intuitively, the most general form of a type is obtained by loosening all inner functional types to relational types, reflecting \cref{rule:castable-graph}.

\begin{definition}[Least general form][lgf][Encoder/Loosening/Castable.html\#SMT.SMTType.lgf]
  \label{def:lgf}
    The least general form of $\alpha : \mathsf{SMTType}$, denoted $\regnotation{\lgf{\alpha}}$ or $\mathsf{lgf}(\alpha)$, is defined structurally on $\alpha$, as follows:
  \begin{align*}
    \lgf{\alpha} &\triangleq \alpha &&\text{if \(\alpha\) is one of } \intS, \boolS, \unitS\\
    \lgf{\alpha} &\triangleq \lgf{\gamma} \toS \optS\ \lgf{\delta} &&\text{if \(\alpha = (\gamma \timesS \delta) \toS \boolS\)}\\
  \end{align*} 
  Remaining cases are defined by structural recursion on $\alpha$.
\end{definition}
%
The intuition behind the least general form of a type is the converse of the most general form: it is obtained by tightening all inner relational types to functional types. Where $\mathsf{mgf}$ \emph{erases} information---forgetting that a relation happens to be a function, as in \cref{rule:castable-graph}---$\mathsf{lgf}$ does the opposite and \emph{invents} it, imposing the determinism and totality that turn a relation back into a function. The two are thus not inverse bijections but a forgetful/free pair, made precise by the Galois connection of the next subsection.

Not only does this provide two normal forms for any type, but it also induces a structure with interesting properties, as depicted in the example below.

\begin{example}
  \label{ex:lgf-mgf}
  Consider the type $\alpha$ defined below:
  \[
  \alpha = \Big((\intS \toS \optS\ \intS) \timesS \big((\intS \timesS \boolS) \toS \boolS\big)\Big) \toS \boolS
  \]
  Then, the least and most general forms of $\alpha$ are:
  \begin{align*}
    \lgf{\alpha} &= (\intS \toS \optS\ \intS) \toS \optS\big(\intS \toS \optS\ \boolS\big)\\
    \mgf{\alpha} &= \Big(\big((\intS \timesS \intS) \toS \boolS\big) \timesS \big((\intS \timesS \boolS) \toS \boolS\big)\Big) \toS \boolS
  \end{align*}
  Moreover, there are several intermediate forms of $\alpha$, defined as follows:
  \begin{align*}
    \alpha^{(1)} &= (\intS \toS \optS\ \intS) \toS \optS\big((\intS \timesS \boolS) \toS \boolS\big)\\
    \alpha^{(2)} &= ((\intS \timesS \intS) \toS \boolS) \toS \optS\big(\intS \toS \optS\ \boolS\big)\\
    \alpha^{(3)} &= ((\intS \timesS \intS) \toS \boolS) \toS \optS\big((\intS \timesS \boolS) \toS \boolS\big)\\
    \alpha^{(4)} &= \Big((\intS \toS \optS\ \intS) \timesS (\intS \toS \optS\ \boolS)\Big) \toS \boolS\\
    \alpha^{(5)} &= \Big(\big((\intS \timesS \intS) \toS \boolS\big) \timesS (\intS \toS \optS\ \boolS)\Big) \toS \boolS
  \end{align*}
  The ordering relationship between the different forms of $\alpha$ is summarized in \Cref{fig:loosening-hasse}, where oriented arrows denote the $\sqsubseteq$ relation:
  \begin{center}
    \begin{tikzpicture}[
      node/.style={draw=none, circle},
      cross/.style={midway, draw=none, fill=green!5, circle, inner sep=0pt, minimum size=8pt},
      node distance=1cm and 1.2cm,
      minimum size=1.5em,
    ]
      % nodes
      \node (bot) at (0,0) {$\lgf{\alpha}$};
      \node[above left=of bot] (a1) {$\alpha^{(1)}$};
      \node[above=of bot] (a4) {$\alpha^{(4)}$};
      \node[above right=of bot] (a2) {$\alpha^{(2)}$};
      \node[above=of a1] (alpha) {$\alpha$};
      \node[above right=of a1] (a3) {$\alpha^{(3)}$};
      \node[above=of a2] (a5) {$\alpha^{(5)}$};
      \node[above=of a3] (top) {$\mgf{\alpha}$};

      % arrows
      \draw[relarr] (bot) -- (a1);
      \draw[relarr] (bot) -- (a2);
      \draw[relarr] (bot) -- (a4);
      \draw[relarr] (a1) to node[cross] {} (a3);
      \draw[relarr] (a1) -- (alpha);
      \draw[relarr] (a2) to node[cross] {} (a3);
      \draw[relarr] (a2) -- (a5);
      \draw[relarr] (a4) to (alpha);
      \draw[relarr] (a4) -- (a5);
      \draw[relarr] (a3) -- (top);
      \draw[relarr] (alpha) -- (top);
      \draw[relarr] (a5) -- (top);

    \end{tikzpicture}
  \end{center}
  \captionof{figure}{The order induced by $\sqsubseteq$ on the eight forms of the type $\alpha$ of this example, from its least general (fully functional) form $\lgf{\alpha}$ at the bottom to its most general (fully relational) form $\mgf{\alpha}$ at the top. Each edge loosens exactly one functional position through \cref{rule:castable-graph}.}%
  \label{fig:loosening-hasse}
\end{example}

\subsection{Galois connection}
\Cref{ex:lgf-mgf} suggests that the structure induced by $\sqsubseteq$ on $\mathsf{SMTType}$ has strong local properties.
We indeed equip the poset $(\mathsf{SMTType},\sqsubseteq)$ with a Galois connection~\cite{davey-priestley} via the adjunction $\mathsf{lgf} \dashv \mathsf{mgf}$.
Although the partially ordered set (poset) is not a lattice globally, each adjoint (closed) interval $[\lgf{\alpha}, \mgf{\alpha}] \coloneq \{ \tau \mid \lgf{\alpha} \sqsubseteq \tau \sqsubseteq \mgf{\alpha} \}$ is a complete lattice.
Every type in $[\lgf{\alpha}, \mgf{\alpha}]$ shares the same least and most general forms, so the interval gathers exactly the mutually loosenable representations of one underlying shape, ordered from the most functional, $\lgf{\alpha}$, to the most relational, $\mgf{\alpha}$; loosening a term always moves upward within a single such interval.
We list below some of the properties of this structure with associated diagrams.
Plain arrows {\tikz[baseline=-.6ex]{\draw[funarr, shorten <=0pt, shorten >=0pt] (0,0) to (1.2em,0);}} denote morphisms of the poset and are always annotated with the corresponding morphism, while dashed gray arrows {\tikz[baseline=-.6ex]{\draw[relarr, shorten <=0pt, shorten >=0pt] (0,0) to (1.2em,0);}} denote the $\sqsubseteq$ relation, read as ``can be loosened into''.


The properties below build up, step by step, to the adjunction $\mathsf{lgf} \dashv \mathsf{mgf}$. First, the two normal forms bracket the type they are computed from.
\begin{theorem}[][lgf\_le, le\_mgf][Encoder/Loosening/Castable.html\#SMT.SMTType.lgf_le, Encoder/Loosening/Castable.html\#SMT.SMTType.le_mgf]
  \(\alpha\) belongs to the interval \([\lgf{\alpha}, \mgf{\alpha}]\), \ie
  \(\mathsf{lgf}(\alpha) \sqsubseteq \alpha \sqsubseteq \mathsf{mgf}(\alpha).\)
  \begin{center}
    \begin{tikzpicture}
      \node[dot] (a) {};
      \node[dot, right=of a] (mgf) {};
      \node[dot, left=of a] (lgf) {};

      \draw[funarr] (a) to [bend left=20] node[above] {$\mathsf{mgf}$} (mgf);
      \draw[funarr] (a) to [bend right=20] node[above] {$\mathsf{lgf}$} (lgf);
      \draw[relarr] (lgf) to [bend right=20] (a);
      \draw[relarr] (a) to [bend right=20] (mgf);
    \end{tikzpicture}
  \end{center}
\end{theorem}
\begin{proof}
By induction on $\alpha$. At atomic types $\lgf{\alpha} = \alpha = \mgf{\alpha}$. At a composite type, each clause of $\mathsf{lgf}$ and $\mathsf{mgf}$ either keeps the head constructor---so the bounds follow from the inductive hypotheses applied componentwise---or trades a functional layer for its graph, a single \ref{rule:castable-graph} step for $\mathsf{mgf}$, and its tightening for $\mathsf{lgf}$.
\end{proof}

Each form is normalizing, or \emph{idempotent}: applying it twice does not change the result.
\begin{theorem}[Idempotence][lgf\_idempotent, mgf\_idempotent][Encoder/Loosening/Castable.html\#SMT.SMTType.lgf_idempotent, Encoder/Loosening/Castable.html\#SMT.SMTType.mgf_idempotent]
  $\mathsf{lgf}$ and $\mathsf{mgf}$ are idempotent, \ie
  \[
    \mathsf{lgf} \circ \mathsf{lgf} = \mathsf{lgf} \quad\text{and}\quad \mathsf{mgf} \circ \mathsf{mgf} = \mathsf{mgf}.
  \]
  \begin{center}
    \begin{tikzpicture}
      \node[dot] (a) {};
      \node[dot, left=of a] (lgf) {};
      \node[dot, right=of a] (mgf) {};

      \draw[funarr] (a) to node[above] {$\mathsf{lgf}$} (lgf);
      \draw[funarr] (lgf) to[loop, in=150, out=210, looseness=30] node[left] {$\mathsf{lgf}$} (lgf);
      \draw[funarr] (a) to node[above] {$\mathsf{mgf}$} (mgf);
      \draw[funarr] (mgf) to[loop, in=30, out=-30, looseness=30] node[right] {$\mathsf{mgf}$} (mgf);
    \end{tikzpicture}
  \end{center}
\end{theorem}
\begin{proof}
By induction on $\alpha$. The type $\mgf{\alpha}$ has no functional layer left for $\mathsf{mgf}$ to rewrite, so a second application is the identity; dually, $\lgf{\alpha}$ has no relational layer for $\mathsf{lgf}$ to tighten.
\end{proof}

Both forms are also monotone for the loosening order \(\castable{}{}\).
\begin{theorem}[Monotonicity][lgf\_mono, mgf\_mono][Encoder/Loosening/Castable.html\#SMT.SMTType.lgf_mono, Encoder/Loosening/Castable.html\#SMT.SMTType.mgf_mono]
  $\mathsf{lgf}$ and $\mathsf{mgf}$ are monotone, \ie if $\alpha \sqsubseteq \beta$, then
  \[
    \mathsf{lgf}(\alpha) \sqsubseteq \mathsf{lgf}(\beta) \quad\text{and}\quad \mathsf{mgf}(\alpha) \sqsubseteq \mathsf{mgf}(\beta)
  \]
  \begin{center}
    \begin{tikzpicture}
      \node[dot] (a) {};
      \node[dot, right=of a] (mgfa) {};
      \node[dot, left=of a] (lgfa) {};
      \node[dot, above= of a] (b) {};
      \node[dot, right=of b] (mgfb) {};
      \node[dot, left=of b] (lgfb) {};
      
      \draw[relarr] (a) to (b);
      \draw[funarr] (a) to node[midway, below] {$\mathsf{mgf}$} (mgfa);
      \draw[funarr] (a) to node[midway, below] {$\mathsf{lgf}$} (lgfa);
      \draw[funarr] (b) to node[midway, above] {$\mathsf{mgf}$} (mgfb);
      \draw[funarr] (b) to node[midway, above] {$\mathsf{lgf}$} (lgfb);
      \draw[relarr] (mgfa) to (mgfb);
      \draw[relarr] (lgfa) to (lgfb);
    \end{tikzpicture}
  \end{center}
\end{theorem}
\begin{proof}
By induction on the derivation of $\castable{\alpha}{\beta}$: the structural rules are carried to derivations between the normal forms via the inductive hypotheses, and \ref{rule:castable-graph} maps to the corresponding step between graphs.
\end{proof}

Monotonicity and idempotence together fix how the two forms compose: each absorbs the other,
a property known as the triangular identities---although we present them both in one diagram here.
\begin{theorem}[Triangular identities][mgf\_of\_lgf, lgf\_of\_mgf][Encoder/Loosening/Castable.html\#SMT.SMTType.mgf_of_lgf, Encoder/Loosening/Castable.html\#SMT.SMTType.lgf_of_mgf]
  $\mathsf{lgf}$ and $\mathsf{mgf}$ satisfy the triangular identities, \ie
  \[
    \mathsf{lgf} \circ \mathsf{mgf} = \mathsf{lgf} \quad\text{and}\quad \mathsf{mgf} \circ \mathsf{lgf} = \mathsf{mgf}
  \]
  \begin{center}
    \begin{tikzpicture}
      \node[dot] (a) {};
      \node[dot, right=of a] (mgfa) {};
      \node[dot, left=of a] (lgfa) {};
      
      \draw[funarr] (a) to node[midway, below] {$\mathsf{mgf}$} (mgfa);
      \draw[funarr] (a) to node[midway, below] {$\mathsf{lgf}$} (lgfa);
      \draw[funarr] (mgfa) to[in=-80, out=-100, looseness=0.7] node[midway, below] {$\mathsf{lgf} \circ \mathsf{mgf}$} (lgfa);
      \draw[funarr] (lgfa) to[in=100, out=80, looseness=0.7] node[midway, above] {$\mathsf{mgf} \circ \mathsf{lgf}$} (mgfa);
    \end{tikzpicture}
  \end{center}
\end{theorem}
\begin{proof}
By induction on $\alpha$; equivalently, $\lgf{\alpha}$ and $\alpha$ lie in the same interval $[\lgf{\alpha}, \mgf{\alpha}]$ and therefore share a most general form, giving $\mathsf{mgf}(\lgf{\alpha}) = \mgf{\alpha}$, and dually $\mathsf{lgf}(\mgf{\alpha}) = \lgf{\alpha}$.
\end{proof}

These facts assemble into the central property of the pair: $\mathsf{lgf}$ and $\mathsf{mgf}$ are adjoint.
\begin{theorem}[Adjunction][castable?.galois][Encoder/Loosening/Castable.html\#SMT.SMTType.castable?.galois]
  $\mathsf{lgf}$ is left adjoint to $\mathsf{mgf}$ on $(\mathsf{SMTType}, \castable{}{})$, denoted $\mathsf{lgf} \dashv \mathsf{mgf}$, i.e., for any $\alpha, \beta : \mathsf{SMTType}$,
  \[
    \mathsf{lgf}(\alpha) \sqsubseteq \beta \iff \alpha \sqsubseteq \mathsf{mgf}(\beta)
  \]
  \begin{center}
    \begin{tikzpicture}
      \node[dot] (a) {};
      \node[dot, left=of a] (lgfa) {};
      \node[dot, above= of a] (b) {};
      \node[dot, right=of b] (mgfb) {};
      
      \draw[funarr] (a) to node[midway, below] {$\mathsf{lgf}$} (lgfa);
      \draw[funarr] (b) to node[midway, above] {$\mathsf{mgf}$} (mgfb);
      \draw[relarr] (lgfa) to coordinate (midl) (b);
      \draw[relarr] (a) to coordinate (midr) (mgfb);
      
      \draw[implies-implies,double equal sign distance, shorten <=5pt, shorten >=5pt, thick] (midl) to (midr);
    \end{tikzpicture}
  \end{center}
\end{theorem}
\begin{proof}
Both directions combine monotonicity, the bounds $\lgf{\alpha} \sqsubseteq \alpha \sqsubseteq \mgf{\alpha}$ and $\lgf{\beta} \sqsubseteq \beta \sqsubseteq \mgf{\beta}$, and the triangular identities. If $\lgf{\alpha} \sqsubseteq \beta$, then $\alpha \sqsubseteq \mgf{\alpha} = \mathsf{mgf}(\lgf{\alpha}) \sqsubseteq \mgf{\beta}$; conversely, if $\alpha \sqsubseteq \mgf{\beta}$, then $\lgf{\alpha} \sqsubseteq \mathsf{lgf}(\mgf{\beta}) = \lgf{\beta} \sqsubseteq \beta$.
\end{proof}

Intuitively, the adjunction means that loosenability can always be decided against a normal form. As $\lgf{\alpha}$ and $\mgf{\beta}$ are the best functional and relational approximations of $\alpha$ and $\beta$, the adjunction says precisely that comparing a type against either canonical approximation suffices.

As mentioned already, meets and joins do not exist globally, but they do exist locally in each adjoint interval $[\lgf{\alpha}, \mgf{\alpha}]$.
They are defined as follows for any $\theta, \tau \in [\lgf{\alpha}, \mgf{\alpha}]$:

\begin{definition}[Join, meet][join, meet][Encoder/Loosening/Castable.html\#SMT.SMTType.join, Encoder/Loosening/Castable.html\#SMT.SMTType.meet]
  The join and meet of $\theta$ and $\tau$ in $[\lgf{\alpha}, \mgf{\alpha}]$, denoted $\theta \regnotation{\sqcup} \tau$ and $\theta \regnotation{\sqcap} \tau$ respectively, are defined as:
  \begin{align*}
    \theta \sqcup \tau &\triangleq \inf\bigl([\theta, \mgf{\theta}] \cap [\tau, \mgf{\tau}]\bigr)\\
    \theta \sqcap \tau &\triangleq \sup\bigl([\lgf{\theta}, \theta] \cap [\lgf{\tau}, \tau]\bigr)
  \end{align*}
\end{definition}

\begin{example}
  Consider the diagram from \Cref{ex:lgf-mgf} again.
  Then, $\alpha^{(1)} \sqcup \alpha^{(2)} = \alpha^{(3)}$ and $\alpha \sqcap \alpha^{(5)} = \alpha^{(4)}$, as depicted below.
  \begin{center}
    \begin{tikzpicture}[
      node/.style={draw=none, circle},
      cross/.style={midway, draw=none, fill=green!5, circle, inner sep=0pt, minimum size=8pt},
      node distance=.8cm and 1cm,
      minimum size=1.5em,
    ]
    \begin{scope}
      % nodes
      \node (bot) at (0,0) {$\lgf{\alpha}$};
      \node[above left=of bot] (a1) {$\alpha^{(1)}$};
      \node[above=of bot] (a4) {$\alpha^{(4)}$};
      \node[above right=of bot] (a2) {$\alpha^{(2)}$};
      \node[above=of a1] (alpha) {$\alpha$};
      \node[above right=of a1, circle, draw=orange, line width=1.2pt, inner sep=1pt] (a3) {$\alpha^{(3)}$};
      \node[above=of a2] (a5) {$\alpha^{(5)}$};
      \node[above=of a3] (top) {$\mgf{\alpha}$};

      % Polygon through (a1, alpha, top, a3), then expanded outward
      \path let
        \p1 = (a1),
        \p2 = (alpha),
        \p3 = (top),
        \p4 = (a3),
        \n1 = {(\x1+\x2+\x3+\x4)/4},
        \n2 = {(\y1+\y2+\y3+\y4)/4},
        \n3 = {atan2(\y1-\n2,\x1-\n1)},
        \n4 = {atan2(\y2-\n2,\x2-\n1)},
        \n5 = {atan2(\y3-\n2,\x3-\n1)},
        \n6 = {atan2(\y4-\n2,\x4-\n1)}
      in
        coordinate (interval-center) at (\n1,\n2)
        coordinate (i-a1) at ($(interval-center)!1.2!(a1.\n3)$)
        coordinate (i-alpha) at ($(interval-center)!1.10!(alpha.\n4)$)
        coordinate (i-top) at ($(interval-center)!1.08!(top.\n5)$)
        coordinate (i-a3) at ($(interval-center)!1.09!(a3.\n6)$);
      % Enforce vertical side edges while keeping top/bottom edges angled
      \path let
        \p1 = (i-a1),
        \p2 = (i-alpha),
        \p3 = (i-top),
        \p4 = (i-a3),
        \p5 = (alpha.north),
        \p6 = (top.north),
        \n1 = {min(\x1,\x2)-1.5},
        \n2 = {max(\x3,\x4)-1.2},
        \n3 = {(\y3-\y2)/(\x3-\x2)},
        \n4 = {max(\y5+2-\n3*\x5,\y6+2-\n3*\x6)},
        \n5 = {\n3*\n1+\n4},
        \n6 = {\n3*\n2+\n4}
      in
        coordinate (i-left-bot) at (\n1,\y1)
        coordinate (i-left-top) at (\n1,\n5)
        coordinate (i-right-top) at (\n2,\n6)
        coordinate (i-right-bot) at (\n2,\y4);
      \draw[draw=thmbar, densely dashed, line width=1pt, rounded corners=10pt]
        (i-left-bot) -- (i-left-top) to node[midway, above, sloped, color=thmbar!80!black, xshift=-1em] {\small$[\alpha^{(1)}, \mgf{\alpha^{(1)}}]$} (i-right-top) -- (i-right-bot) -- cycle;

      % Polygon through (a3, top, a5, a2), then expanded outward
      \path let
        \p1 = (a3),
        \p2 = (top),
        \p3 = (a5),
        \p4 = (a2),
        \n1 = {(\x1+\x2+\x3+\x4)/4},
        \n2 = {(\y1+\y2+\y3+\y4)/4},
        \n3 = {atan2(\y1-\n2,\x1-\n1)},
        \n4 = {atan2(\y2-\n2,\x2-\n1)},
        \n5 = {atan2(\y3-\n2,\x3-\n1)},
        \n6 = {atan2(\y4-\n2,\x4-\n1)}
      in
        coordinate (j-center) at (\n1,\n2)
        coordinate (j-a3) at ($(j-center)!1.09!(a3.\n3)$)
        coordinate (j-top) at ($(j-center)!1.08!(top.\n4)$)
        coordinate (j-a5) at ($(j-center)!1.10!(a5.\n5)$)
        coordinate (j-a2) at ($(j-center)!1.16!(a2.\n6)$);
      % Enforce vertical side edges while keeping top/bottom edges angled
      \path let
        \p1 = (j-a3),
        \p2 = (j-top),
        \p3 = (j-a5),
        \p4 = (j-a2),
        \p5 = (top.north),
        \p6 = (a5.north),
        \n1 = {min(\x1,\x2)+0.4},
        \n2 = {max(\x3,\x4)+1.1},
        \n3 = {(\y4-\y1)/(\x4-\x1)},
        \n4 = {max(\y5+2-\n3*\x5,\y6+2-\n3*\x6)},
        \n5 = {\n3*\n1+\n4},
        \n6 = {\n3*\n2+\n4}
      in
        coordinate (j-left-bot) at (\n1,\y1-2mm)
        coordinate (j-left-top) at (\n1,\n5)
        coordinate (j-right-top) at (\n2,\n6)
        coordinate (j-right-bot) at (\n2,\y4-4mm);
      \draw[draw=rmkbar, densely dashed, line width=1pt, rounded corners=10pt]
        (j-left-bot) -- (j-left-top) to node[midway, above, sloped, color=rmkbar!80!black, xshift=1em] {\small$[\alpha^{(2)}, \mgf{\alpha^{(2)}}]$} (j-right-top) -- (j-right-bot) -- cycle;

      % arrows
      \draw[relarr] (bot) -- (a1);
      \draw[relarr] (bot) -- (a2);
      \draw[relarr] (bot) -- (a4);
      \draw[relarr] (a1) to node[cross] {} (a3);
      \draw[relarr] (a1) -- (alpha);
      \draw[relarr] (a2) to node[cross] {} (a3);
      \draw[relarr] (a2) -- (a5);
      \draw[relarr] (a4) to (alpha);
      \draw[relarr] (a4) -- (a5);
      \draw[relarr] (a3) -- (top);
      \draw[relarr] (alpha) -- (top);
      \draw[relarr] (a5) -- (top);
    \end{scope}
    \begin{scope}[xshift=6.5cm]
      % nodes
      \node (bot) at (0,0) {$\lgf{\alpha}$};
      \node[above left=of bot] (a1) {$\alpha^{(1)}$};
      \node[above=of bot, circle, draw=orange, line width=1.2pt, inner sep=1pt] (a4) {$\alpha^{(4)}$};
      \node[above right=of bot] (a2) {$\alpha^{(2)}$};
      \node[above=of a1] (alpha) {$\alpha$};
      \node[above right=of a1] (a3) {$\alpha^{(3)}$};
      \node[above=of a2] (a5) {$\alpha^{(5)}$};
      \node[above=of a3] (top) {$\mgf{\alpha}$};

      % Frame for [\lgf{\alpha}, \alpha] through (a1, alpha, a4, bot)
      \path let
        \p1 = (a1),
        \p2 = (alpha),
        \p3 = (a4),
        \p4 = (bot),
        \n1 = {(\x1+\x2+\x3+\x4)/4},
        \n2 = {(\y1+\y2+\y3+\y4)/4},
        \n3 = {atan2(\y1-\n2,\x1-\n1)},
        \n4 = {atan2(\y2-\n2,\x2-\n1)},
        \n5 = {atan2(\y3-\n2,\x3-\n1)},
        \n6 = {atan2(\y4-\n2,\x4-\n1)}
      in
        coordinate (k-center) at (\n1,\n2)
        coordinate (k-a1) at ($(k-center)!1.12!(a1.\n3)$)
        coordinate (k-alpha) at ($(k-center)!1.08!(alpha.\n4)$)
        coordinate (k-a4) at ($(k-center)!1.08!(a4.\n5)$)
        coordinate (k-bot) at ($(k-center)!1.12!(bot.\n6)$);
      \path let
        \p1 = (k-a1),
        \p2 = (k-alpha),
        \p3 = (k-a4),
        \p4 = (k-bot),
        \p5 = (alpha.north),
        \p6 = (a4.north),
        \n1 = {min(\x1,\x2)-0.7},
        \n2 = {max(\x3,\x4)+0.9},
        \n3 = {(\y3-\y2)/(\x3-\x2)},
        \n4 = {max(\y5+1.8-\n3*\x5,\y6+1.8-\n3*\x6)},
        \n5 = {\n3*\n1+\n4},
        \n6 = {\n3*\n2+\n4}
      in
        coordinate (k-left-bot) at (\n1,\y1)
        coordinate (k-left-top) at (\n1,\n5)
        coordinate (k-right-top) at (\n2,\n6)
        coordinate (k-right-bot) at (\n2,\y4-4mm);
      \draw[draw=thmbar, densely dashed, line width=1pt, rounded corners=10pt]
        (k-left-bot) -- (k-left-top) -- (k-right-top) -- (k-right-bot) to node[midway, below, sloped, color=thmbar!80!black, xshift=-1em] {$[\lgf{\alpha}, \alpha]$} cycle;

      % Frame for [\lgf{\alpha}, \alpha^{(5)}] through (a4, a5, a2, bot)
      \path let
        \p1 = (a4),
        \p2 = (a5),
        \p3 = (a2),
        \p4 = (bot),
        \n1 = {(\x1+\x2+\x3+\x4)/4},
        \n2 = {(\y1+\y2+\y3+\y4)/4},
        \n3 = {atan2(\y1-\n2,\x1-\n1)},
        \n4 = {atan2(\y2-\n2,\x2-\n1)},
        \n5 = {atan2(\y3-\n2,\x3-\n1)},
        \n6 = {atan2(\y4-\n2,\x4-\n1)}
      in
        coordinate (m-center) at (\n1,\n2)
        coordinate (m-a4) at ($(m-center)!1.08!(a4.\n3)$)
        coordinate (m-a5) at ($(m-center)!1.08!(a5.\n4)$)
        coordinate (m-a2) at ($(m-center)!1.12!(a2.\n5)$)
        coordinate (m-bot) at ($(m-center)!1.12!(bot.\n6)$);
      \path let
        \p1 = (m-a4),
        \p2 = (m-a5),
        \p3 = (m-a2),
        \p4 = (m-bot),
        \p5 = (a5.north),
        \p6 = (a4.north),
        \n1 = {min(\x1,\x4)-0.9},
        \n2 = {max(\x2,\x3)+0.9},
        \n3 = {(\y2-\y1)/(\x2-\x1)},
        \n4 = {max(\y5+1.8-\n3*\x5,\y6+1.8-\n3*\x6)},
        \n5 = {\n3*\n1+\n4},
        \n6 = {\n3*\n2+\n4}
      in
        coordinate (m-left-bot) at (\n1,\y4-4mm)
        coordinate (m-left-top) at (\n1,\n5)
        coordinate (m-right-top) at (\n2,\n6)
        coordinate (m-right-bot) at (\n2,\y3);
      \draw[draw=rmkbar, densely dashed, line width=1pt, rounded corners=10pt]
        (m-left-bot) -- (m-left-top) -- (m-right-top) -- (m-right-bot) to node[midway, below, sloped, color=rmkbar!80!black, xshift=1em] {$[\lgf{\alpha}, \alpha^{(5)}]$} cycle;

      % arrows
      \draw[relarr] (bot) -- (a1);
      \draw[relarr] (bot) -- (a2);
      \draw[relarr] (bot) -- (a4);
      \draw[relarr] (a1) to node[cross] {} (a3);
      \draw[relarr] (a1) -- (alpha);
      \draw[relarr] (a2) to node[cross] {} (a3);
      \draw[relarr] (a2) -- (a5);
      \draw[relarr] (a4) to (alpha);
      \draw[relarr] (a4) -- (a5);
      \draw[relarr] (a3) -- (top);
      \draw[relarr] (alpha) -- (top);
      \draw[relarr] (a5) -- (top);
    \end{scope}
    \end{tikzpicture}
  \end{center}
  Each is computed, inside the lattice of \Cref{fig:loosening-hasse}, as an extremum of an intersection of adjoint intervals (framed): on the left, the join $\alpha^{(1)} \sqcup \alpha^{(2)} = \alpha^{(3)}$ is the least element of $[\alpha^{(1)},\mgf{\alpha^{(1)}}] \cap [\alpha^{(2)},\mgf{\alpha^{(2)}}]$; on the right, the meet $\alpha \sqcap \alpha^{(5)} = \alpha^{(4)}$ is the greatest element of $[\lgf{\alpha},\alpha] \cap [\lgf{\alpha},\alpha^{(5)}]$.
\end{example}

\begin{note}
The next four order-theoretic results are the only statements in this chapter that are not formalized in Lean. They are not used by the certified encoder; we include short mathematical proofs for completeness. As the short proofs below show, mechanizing them would be a routine exercise.
\end{note}

We then show expected properties of the join and meet operations, such as idempotence, commutativity, and absorption.

\begin{lemma}[Idempotence]
  For any SMT type $\theta$, $\theta \sqcup \theta = \theta$ and $\theta \sqcap \theta = \theta$.
\end{lemma}
\begin{proof}[Proof sketch]
Immediate from the definitions: $\theta \sqcup \theta = \inf([\theta, \mgf{\theta}]) = \theta$ and $\theta \sqcap \theta = \sup([\lgf{\theta}, \theta]) = \theta$.
\end{proof}

\begin{lemma}[Commutativity]
  For any SMT types $\theta$ and $\tau$, $\theta \sqcup \tau = \tau \sqcup \theta$ and $\theta \sqcap \tau = \tau \sqcap \theta$.
\end{lemma}
\begin{proof}[Proof sketch]
The defining intersections are symmetric in $\theta$ and $\tau$, so both operations are commutative.
\end{proof}

\begin{lemma}[Absorption]
  For any SMT types $\theta$ and $\tau$, $\theta \sqcup (\theta \sqcap \tau) = \theta$ and $\theta \sqcap (\theta \sqcup \tau) = \theta$.
\end{lemma}
\begin{proof}[Proof sketch]
Since $\theta \sqcap \tau \sqsubseteq \theta$, the join $\theta \sqcup (\theta \sqcap \tau)$ is the supremum of two elements bounded above by $\theta$, hence $\theta$; dually, $\theta \sqsubseteq \theta \sqcup \tau$ gives $\theta \sqcap (\theta \sqcup \tau) = \theta$.
\end{proof}

\begin{theorem}[Complete lattice]
  \label{thm:loosening-lattice}
  For any SMT type $\alpha$, the adjoint interval $[\lgf{\alpha}, \mgf{\alpha}]$ is a complete lattice with join $\sqcup$, meet $\sqcap$, bottom $\lgf{\alpha}$ and top $\mgf{\alpha}$.
\end{theorem}
\begin{proof}[Proof sketch]
The interval is finite: each of the finitely many functional positions of $\alpha$ makes an independent choice between its functional and relational forms. Because the interval is finite and closed under binary joins and meets, every finite subset has a join and a meet by iteration, using the bottom and top elements for the empty subset. Every subset of a finite interval is finite; hence arbitrary joins and meets exist, so the lattice is complete.
\end{proof}

\begin{remark}[Loosening is not ordinary subtyping]
Although $\sqsubseteq$ is a partial order on SMT types, it should not be confused with ordinary subtyping. A judgment $\castable{\alpha}{\beta}$ does not license subsumption: it records an explicit change of SMT representation that preserves the underlying B set. Function types make the distinction clear. Ordinary function subtyping is contravariant in the domain and covariant in the codomain, whereas rule \ref{rule:castable-fun} loosens both components covariantly, extending the transported function by a default value outside the image of the domain cast. Loosening is therefore a semantics-preserving representation change, rather than a subtyping discipline internal to SMT-LIB.
\end{remark}

\section{Loosening equations}
\label{sec:loosening-rules}

Having settled when and how a \emph{type} may be loosened, we now define how \emph{terms} are loosened accordingly. The following construction $\regnotation{\mathsf{loosen}_{\mathbb{c}}}$ takes a concrete SMT term $x$---expected to be of SMT type $\alpha$---and a cast path $\mathbb{c} : \castpath{\alpha}{\beta}$, and returns a pair
\[
  \langle \regnotation{\lvar{x}}, \regnotation{\spec{\lvar{x}}} \rangle := \mathsf{loosen}_{\mathbb{c}}(x)
\]
made of a \emph{fresh} variable $\lvar{x}$ of the target type $\beta$ and a boolean term $\spec{\lvar{x}}$, its \emph{specification}, which ties $\lvar{x}$ to $x$. Informally, $\spec{\lvar{x}}$ asserts that ``$\lvar{x}$ denotes the image of $x$ under the cast encoded by $\mathbb{c}$''; \Cref{sec:loosening-correctness} turns this slogan into a theorem.

The definition proceeds by structural recursion on $\mathbb{c}$, with one equation per cast-path constructor, each represented in the following \Cref{case:refl,case:opt,case:pair,case:chpred,case:fun,case:graph}.
We write source variables in gray, as $\var{z}$, and their loosened counterparts with a bang, as $\lvar{z}$; we abbreviate a recursive call $\langle \lvar{z}, \spec{\lvar{z}} \rangle := \mathsf{loosen}_{\mathbb{c}'}(\var{z})$ by saying that $\lvar{z}$ is \emph{obtained by loosening} $\var{z}$, thereby also obtaining its specification $\spec{\lvar{z}}$.
Each case introduces the fresh (loosened) target variable and also generates a specification that relates it to the source term.
The equations below are the proof-oriented formulation; the implementation additionally short-circuits reflexive sub-casts (replacing a recursive call on $\castpath{\gamma}{\gamma}$ by an identity), an optimization that affects neither the denotation nor the correctness argument and that we therefore suppress.

\subsection{Case \texttt{refl}}\label{case:refl}
Every atomic type is loosened to itself, and the specification merely equates the fresh variable with the source term.
\begin{definition}[refl][loosen][Encoder/Loosening/Rules.html\#loosenAux_prf]
Let $\mathbb{c} = \mathtt{refl} : \castpath{\alpha}{\alpha}$---which exists exactly when $\alpha$ is one of $\intS$, $\boolS$, and $\unitS$---and let $\lvar{x}$ be fresh. Define $\spec{\lvar{x}}$ by
\[
  \spec{\lvar{x}} \;\;\triangleq\;\; \lvar{x} \eqS x
\]
\end{definition}

\subsection{Case \texttt{opt}}\label{case:opt}
An optional value is loosened by cases on its two constructors. When the source term is syntactically $\noneS$ or $\someS$, we commit to the matching target shape directly; otherwise we emit a dynamic test on $x$.

\begin{definition}[opt][loosen][Encoder/Loosening/Rules.html\#loosenAux_prf]
Let $\mathbb{c} = \mathtt{opt}\;\mathbb{c}_\alpha : \castpath{\optS\,\alpha}{\optS\,\alpha'}$, where $\mathbb{c}_\alpha : \castpath{\alpha}{\alpha'}$, and let $\lvar{x}$ be fresh. Define $\spec{\lvar{x}}$ by
\[
\begin{array}{r@{\qquad}rl}
x = \noneS_{\alpha}: & \spec{\lvar{x}} &\triangleq \lvar{x} \eqS \noneS_{\alpha'}\\
x = \someS\ y : & \spec{\lvar{x}} &\triangleq \existsS \lvar{y} : \alpha'\cdot\, (\lvar{x} \eqS \someS\ \lvar{y}) \andS \spec{\lvar{y}} \\[1ex]
& & \text{where } \langle \lvar{y}, \spec{\lvar{y}} \rangle \text{ is obtained by loosening } y\\[1ex]
\text{otherwise}: & \spec{\lvar{x}} &\triangleq
\begin{aligned}[t]
&\iteS (x \eqS \noneS_{\alpha})\, (\lvar{x} \eqS \noneS_{\alpha'})\\
&\quad \left(\existsS \lvar{y} : \alpha'\cdot\, (\lvar{x} \eqS \someS\ \lvar{y}) \andS \spec{\lvar{y}}\right)
\end{aligned} \\
& & \text{where } \langle \lvar{y}, \spec{\lvar{y}} \rangle \text{ is obtained by loosening } \theS x
\end{array}
\]
\end{definition}


\subsection{Case \texttt{pair}}\label{case:pair}
A pair is loosened componentwise: the two projections are loosened independently and reassembled.

\begin{definition}[pair][loosen][Encoder/Loosening/Rules.html\#loosenAux_prf]
Let $\mathbb{c} = \mathtt{pair}\;\mathbb{c}_\alpha\;\mathbb{c}_\beta : \castpath{\alpha \timesS \beta}{\alpha' \timesS \beta'}$, where $\mathbb{c}_\alpha : \castpath{\alpha}{\alpha'}$ and $\mathbb{c}_\beta : \castpath{\beta}{\beta'}$, and let $\lvar{x}$ be fresh. Define $\spec{\lvar{x}}$ by
\[
\spec{\lvar{x}} \;\;\triangleq\;\; \existsS \lvar{fst} : \alpha',\, \lvar{snd} : \beta'\cdot\, (\lvar{x} \eqS \pairS{\lvar{fst}}{\lvar{snd}}) \andS \spec{\lvar{fst}} \andS \spec{\lvar{snd}}
\]
where $\langle \lvar{fst}, \spec{\lvar{fst}} \rangle$ and $\langle \lvar{snd}, \spec{\lvar{snd}} \rangle$ are obtained by loosening $\fstS{x}$ and $\sndS{x}$ respectively.
\end{definition}

\subsection{Case \texttt{chpred}}\label{case:chpred}
A characteristic predicate is loosened by relabeling its domain along the inner cast: the loosened target $\lvar{x}$ must hold for an element $\lvar{z}$ loosening $\var{z}$ exactly when $x$ holds for that source element $\var{z}$.

\begin{definition}[chpred][loosen][Encoder/Loosening/Rules.html\#loosenAux_prf]
Let $\mathbb{c} = \mathtt{chpred}\;\mathbb{c}_\alpha : \castpath{(\alpha \toS \boolS)}{(\alpha' \toS \boolS)}$, where $\mathbb{c}_\alpha : \castpath{\alpha}{\alpha'}$, and let $\lvar{x}$ be fresh. Let $\var{z}:\alpha$ be fresh and set
\[
  \langle\lvar{z},\spec{\lvar{z}}\rangle
  \coloneq \mathsf{loosen}_{\mathbb{c}_\alpha}(\var{z}).
\]
Define $\spec{\lvar{x}}$ by
\[
\spec{\lvar{x}} \;\;\triangleq\;\; \lvar{x} \eqS \lamS \lvar{z} : \alpha' \cdot\, \existsS \var{z} : \alpha\cdot\, \appS x (\var{z}) \andS \spec{\lvar{z}}.
\]
The recursive specification $\spec{\lvar{z}}$ relates $\var{z}$ to its image at type $\alpha'$.
\end{definition}

\subsection{Case \texttt{fun}}\label{case:fun}
When the codomain is not $\boolS$---so the source is a genuine function rather than a characteristic predicate---loosening acts pointwise: the loosened function $\lvar{x}$ agrees with $x$ on every argument that lies in the image of the domain cast, and takes a fixed \emph{default} value elsewhere. Pinning down the value off the image is what makes the loosened function total, and hence well typed; we therefore first specify these defaults.

We write $\regnotation{\mathsf{defaultSpec}_{\gamma}}(t)$ for the predicate stating
that $t$ is the canonical default element of type $\gamma$, defined recursively by
\[
\mathsf{defaultSpec}_{\gamma}(t) \triangleq
\begin{cases}
t \eqS 0 & \text{if } \gamma = \intS\\
t \eqS \mathsf{false} & \text{if } \gamma = \boolS\\
\mathsf{true} & \text{if } \gamma = \unitS\\
t \eqS \noneS_{\gamma_1} & \text{if } \gamma = \optS\ \gamma_1\\
\begin{aligned}[t]
&\mathsf{defaultSpec}_{\gamma_1}(\fstS t) \andS \mathsf{defaultSpec}_{\gamma_2}(\sndS t)
\end{aligned} & \text{if } \gamma = \gamma_1 \timesS \gamma_2\\
\forallS \var{u} : \gamma_1\cdot\, \mathsf{defaultSpec}_{\gamma_2}(\appS t (\var{u})) & \text{if } \gamma = \gamma_1 \toS \gamma_2
\end{cases}
\]

\begin{example}
  Let $\gamma \coloneq (\intS \timesS \intS) \toS \intS$. The default specification, at type $\gamma$, of the term \(\lamS \var{x} : \intS \timesS \intS \cdot \fstS \var{x} +  \sndS \var{x}\) is given by the following SMT term:
  \[
  \begin{aligned}
  &\mathsf{defaultSpec}_{\gamma}\left(\lamS \var{x} : \intS \timesS \intS \cdot \fstS \var{x} +  \sndS \var{x}\right) = \\
  &\mkern50mu \forallS \var{u} : \intS \timesS \intS\cdot\, \mathsf{defaultSpec}_{\intS} \left( \appS{\left(\lamS \var{x} : \intS \timesS \intS \cdot \fstS \var{x} +  \sndS \var{x}\right)}{\var{u}}\right) =\\
  &\mkern50mu \forallS \var{u} : \intS \timesS \intS\cdot\, \left(\appS{\left(\lamS \var{x} : \intS \timesS \intS \cdot \fstS \var{x} +  \sndS \var{x}\right)}{\var{u}}\right) \eqS 0
  \end{aligned}
  \]

  Note that this term could be $\beta$-reduced to \(\forallS \var{u} : \intS \timesS \intS\cdot\, \fstS \var{u} +  \sndS \var{u} \eqS 0\), but we do not simplify terms at this stage.
\end{example}

\begin{definition}[fun][loosen][Encoder/Loosening/Rules.html\#loosenAux_prf]
Let $\mathbb{c} = \mathtt{fun}\;\mathbb{c}_\alpha\;\mathbb{c}_\beta : \castpath{(\alpha \toS \beta)}{(\alpha' \toS \beta')}$, where $\mathbb{c}_\alpha : \castpath{\alpha}{\alpha'}$, $\mathbb{c}_\beta : \castpath{\beta}{\beta'}$, and $\beta \neq \boolS$. Let $\lvar{x}$ be fresh. Define $\spec{\lvar{x}}$ by
  \vspace{-.5\baselineskip}
\begin{align*}
  \spec{\lvar{x}} \;\;\triangleq\;\; {} & \forallS \lvar{a} : \alpha'\cdot\, \\
  &\qquad \iteS \left(\existsS \var{a} : \alpha\cdot\, \spec{\lvar{a}}\right) \\
  &\qquad\quad \left(
    \begin{aligned}
      &\forallS \lvar{b} : \beta'\cdot\, ((\appS{\lvar{x}}{(\lvar{a})} \eqS \lvar{b}) \eqS {}\\
      &\qquad (\existsS \var{a} : \alpha,\, \var{b} : \beta\cdot\, (\appS x (\var{a}) \eqS \var{b}) \andS (\spec{\lvar{a}} \andS \spec{\lvar{b}})))
    \end{aligned}
  \right) \\
  &\qquad\quad \left(\mathsf{defaultSpec}_{\beta'}(\appS{\lvar{x}}{(\lvar{a})})\right)
\end{align*}
where $\langle \lvar{a}, \spec{\lvar{a}} \rangle$ and $\langle \lvar{b}, \spec{\lvar{b}} \rangle$ are obtained by loosening fresh source variables $\var{a} : \alpha$ and $\var{b} : \beta$ respectively.
Thus, at an argument in the image of the domain cast, the specification fixes the result to the cast of the source-function result; off that image, it fixes the result to the canonical default of $\beta'$.
\end{definition}

\subsection{Case \texttt{graph}}\label{case:graph}
This case turns a partial function into its graph: the functional encoding $\alpha \toS \optS\,\beta$ is loosened to the relational encoding $(\alpha' \timesS \beta') \toS \boolS$. The fresh variable $\lvar{x}$ is the characteristic predicate of the set of loosened pairs $\lvar{z}$ such that $x$ maps the first component of $\var{z}$ to its second.

\begin{definition}[graph][loosen][Encoder/Loosening/Rules.html\#loosenAux_prf]
Let $\mathbb{c} = \mathtt{graph}\;\mathbb{c}_\alpha\;\mathbb{c}_\beta : \castpath{(\alpha \toS \optS\,\beta)}{((\alpha' \timesS \beta') \toS \boolS)}$, where $\mathbb{c}_\alpha : \castpath{\alpha}{\alpha'}$ and $\mathbb{c}_\beta : \castpath{\beta}{\beta'}$, and let $\lvar{x}$ be fresh. Let
\[
  \mathbb{c}_{\alpha \timesS \beta} \coloneq (\mathbb{c}_\alpha,\mathbb{c}_\beta)
  : \castpath{\alpha \timesS \beta}{\alpha' \timesS \beta'},
\]
and, for a fresh source pair $\var{z}:\alpha \timesS \beta$, set
\[
  \langle \lvar{z}, \spec{\lvar{z}} \rangle
  \coloneq \mathsf{loosen}_{\mathbb{c}_{\alpha \timesS \beta}}(\var{z}).
\]
Define $\spec{\lvar{x}}$ by
  \vspace{-.5\baselineskip}
\begin{align*}
\spec{\lvar{x}} \;\;\triangleq\;\; {} & \lvar{x} \eqS \lamS \lvar{z} : \alpha' \timesS \beta' \cdot\\
&\qquad \existsS \var{z} : \alpha \timesS \beta\cdot\, (\appS x (\fstS \var{z}) \eqS \someS(\sndS \var{z})) \andS \spec{\lvar{z}}
\end{align*}
\end{definition}

\section{Correctness of loosening}
\label{sec:loosening-correctness}

It remains to discharge the slogan of \Cref{sec:loosening-rules}, that $\spec{\lvar{x}}$ pins $\lvar{x}$ to the cast image of $x$. We make this precise semantically, against the ZFC interpretation $\interpZ{\cdot}$ of \Cref{sec:smt-semantics}.

\subsection{Semantic casts}
\label{subsec:semantic-casts}
Every cast path $\mathbb{c} : \castpath{\alpha}{\beta}$ induces a \emph{semantic cast}: a ZFC function $\zop{\mathbb{c}} \in {\interpZ{\beta}}^{\interpZ{\alpha}}$ that sends each constructor to the canonical set-theoretic coercion between the interpretations of its source and target. Formally available under the name {\small\leancodeptr{castZF\_of\_path}{SMT/Reasoning/LooseningDefs.html\#castZF_of_path}}, and defined by recursion on $\mathbb{c}$, it is the semantic shadow of $\mathsf{loosen}_{\mathbb{c}}$, in that where the construction rewrites a \emph{term}, the semantic cast transports the \emph{set} that term denotes. We write the casts as $\lamZ$-abstractions in the notation of \Cref{subsec:zflean-func-eval-lambda}, and record with each clause the set-typing $\zop{\mathbb{c}} \in {\interpZ{\beta}}^{\interpZ{\alpha}}$ of the cast it defines, since the component casts reappear on the right-hand sides of later clauses.

The first five clauses are \emph{structural}: they loosen componentwise and leave the shape of the denoted value untouched, mirroring the homomorphic rules of $\sqsubseteq$. Only \ref{rule:castable-graph} changes shape, trading the functional encoding $\alpha \toS \optS\,\beta$ of a partial function for the relational encoding $(\alpha' \timesS \beta') \toS \boolS$ of its graph; its set-theoretic content is the only one not immediate from the constructor.

\paragraph{\texttt{refl}.} A reflexive path acts as the identity, \ie from \(\mathbb{c} : \castpath{\alpha}{\alpha}\),
\begin{equation}
  \label{eq:castZF-refl}
  \zop{\mathbb{c}} \triangleq \idZ{\interpZ{\alpha}} \;\in\; {\interpZ{\alpha}}^{\interpZ{\alpha}}. \tag{$\zop{\mathtt{refl}}$}
\end{equation}

\paragraph{\texttt{opt}.} The option cast loosens the value under $\someS$ constructor-wise; from \(\mathbb{c} : \castpath{\optS\ \alpha}{\optS\ \alpha'}\), we obtain \(\mathbb{c}_\alpha : \castpath{\alpha}{\alpha'}\) and define \(\zop{\mathbb{c}} \in {\interpZ{\optS\,\alpha'}}^{\interpZ{\optS\,\alpha}}\) as
\begin{equation}
  \label{eq:castZF-opt}
\zop{\mathbb{c}} \;\triangleq\;
\begin{cases}
  \;\noneZ &\;\mapsto\; \noneZ\\
  \;\someZ\ a &\;\mapsto\; \someZ\ \appZ \zop{\mathbb{c}}_\alpha(a).
\end{cases}
\tag{$\zop{\mathtt{opt}}$}
\end{equation}

\paragraph{\texttt{pair}.} The pair cast is the product of the component casts; from \(\mathbb{c} : \castpath{\alpha \timesS \beta}{\alpha' \timesS \beta'}\), we obtain \(\mathbb{c}_\alpha : \castpath{\alpha}{\alpha'}\) and \(\mathbb{c}_\beta : \castpath{\beta}{\beta'}\) and define \(\zop{\mathbb{c}} \in {\interpZ{\alpha' \timesS \beta'}}^{\interpZ{\alpha \timesS \beta}}\) as
\begin{equation}
  \label{eq:castZF-pair}
  \zop{\mathbb{c}} \;\triangleq\;
  \lamZ\,(a, b) \mapsto \bigl(\appZ \zop{\mathbb{c}}_\alpha(a),\; \appZ \zop{\mathbb{c}}_\beta(b)\bigr).
  \tag{$\zop{\mathtt{pair}}$}
\end{equation}

\paragraph{\texttt{chpred}.} The characteristic-predicate cast transports predicates along casts; from \(\mathbb{c} : \castpath{(\alpha \toS \boolS)}{(\alpha' \toS \boolS)}\), we obtain \(\mathbb{c}_\alpha : \castpath{\alpha}{\alpha'}\) and define \(\zop{\mathbb{c}} \in {\interpZ{\alpha' \toS \boolS}}^{\interpZ{\alpha \toS \boolS}}\) as
\begin{equation}
  \label{eq:castZF-chpred}
  \zop{\mathbb{c}} \;\triangleq\;
  \lamZ\,X \mapsto \left(\lamZ\,a' \mapsto
  \begin{cases}
    \appZ X\bigl(\appZ {\zop{\mathbb{c}}}_\alpha^{-1}(a')\bigr) & \text{if } a' \in \Ran(\zop{\mathbb{c}}_\alpha),\\
    \botZ & \text{otherwise}
  \end{cases}\right).
  \tag{$\zop{\mathtt{chpred}}$}
\end{equation}

\paragraph{\texttt{fun}.} The function cast conjugates a function by the component casts; from \(\mathbb{c} : \castpath{(\alpha \toS \beta)}{(\alpha' \toS \beta')}\) with $\beta \neq \boolS$, we obtain \(\mathbb{c}_\alpha : \castpath{\alpha}{\alpha'}\) and \(\mathbb{c}_\beta : \castpath{\beta}{\beta'}\) and define \(\zop{\mathbb{c}} \in {\interpZ{\alpha' \toS \beta'}}^{\interpZ{\alpha \toS \beta}}\) as
\begin{equation}
  \label{eq:castZF-fun}
  \zop{\mathbb{c}} \;\triangleq\;
  \lamZ\,X \mapsto \left(\lamZ\,a' \mapsto
  \begin{cases}
    \appZ {\zop{\mathbb{c}}}_\beta\bigl(\appZ X(\appZ {\zop{\mathbb{c}}}_\alpha^{-1}(a'))\bigr) & \text{if } a' \in \Ran(\zop{\mathbb{c}}_\alpha),\\
    d_{\beta'} & \text{otherwise}
  \end{cases}\right).
  \tag{$\zop{\mathtt{fun}}$}
\end{equation}
where $d_{\beta'}$ is defined by
\[
  \interpS{\mathsf{default}(\beta')} = \denval{d_{\beta'}}[\beta'],
\]
so $d_{\beta'} \in \interpZ{\beta'}$ is the canonical default value (\Cref{def:smt-default}).

\paragraph{\texttt{graph}.} The graph cast carries a partial function $F \in (\optZ\,\interpZ{\beta})^{\interpZ{\alpha}}$ to the characteristic predicate of its graph, pushed forward through the two component casts; from \(\mathbb{c} : \castpath{(\alpha \toS \optS\,\beta)}{((\alpha' \timesS \beta') \toS \boolS)}\), we obtain \(\mathbb{c}_\alpha : \castpath{\alpha}{\alpha'}\) and \(\mathbb{c}_\beta : \castpath{\beta}{\beta'}\) and define \(\zop{\mathbb{c}} \in {\interpZ{(\alpha' \timesS \beta') \toS \boolS}}^{\interpZ{\alpha \toS \optS\,\beta}}\) as
\begin{equation}
  \label{eq:castZF-graph}
  \zop{\mathbb{c}} \;\triangleq\;
  \lamZ\,F \mapsto \left(\lamZ\,(a',b') \mapsto
  \begin{cases}
    \topZ & \begin{aligned}[t]&\text{if } a' \in \Ran(\zop{\mathbb{c}}_\alpha),\ b' \in \Ran(\zop{\mathbb{c}}_\beta),\\ &\text{and } \appZ F(\appZ{\zop{\mathbb{c}}}_\alpha^{-1}(a')) = \someZ\ \appZ{\zop{\mathbb{c}}}_\beta^{-1}(b'),\end{aligned}\\
    \botZ & \text{otherwise}
  \end{cases}\right).
\tag{$\zop{\mathtt{graph}}$}
\end{equation}

Each $\zop{\mathbb{c}}$ is injective---the formal counterpart of the statement that loosening changes the \emph{representation} of a value, never the value itself---and reflects the expected behavior of a cast.

\begin{lemma}[Injectivity of semantic casts][castZF\_of\_path\_injective][SMT/Reasoning/LooseningDefs.html\#castZF_of_path_injective]
\label{lem:semantic-cast-injective}
For every cast path $\mathbb{c} : \castpath{\alpha}{\beta}$, the semantic cast $\zop{\mathbb{c}}$ is injective.
\end{lemma}

\begin{proof}
By induction on $\mathbb{c}$. The \texttt{refl} case is the identity function, which is injective. The structural cases preserve injectivity componentwise; in the \texttt{chpred}, \texttt{fun}, and \texttt{graph} cases the pre-image reconstructions used above are well-defined precisely because the component casts are injective by the inductive hypothesis.
\end{proof}

\subsection{Correctness statement}

Given $\stype{\Gamma}{x}{\alpha}$ and a fresh variable $\lvar{x}$, we abbreviate, as before,
\[
\langle \lvar{x}, \spec{\lvar{x}} \rangle := \mathsf{loosen}_{\mathbb{c}}(x).
\]
The correctness statement proved for the loosening construction is the following.

\begin{theorem}[Correctness of loosening][loosenAux\_prf\_exact][SMT/Reasoning/Basic/LoosenAuxExact.html\#loosenAux_prf_exact]
\label{thm:loosening-correctness}
Let $\Gamma$ be a type-context, $x$ be a concrete term, and $\alpha$ be an SMT type such that $\stype{\Gamma}{x}{\alpha}$.
We also assume that we have a cast path $\mathbb{c} : \castpath{\alpha}{\beta}$ for an SMT type $\beta$.

Let $\langle \lvar{x}, \spec{\lvar{x}} \rangle \coloneq \mathsf{loosen}_{\mathbb{c}}(x)$.
Then:
\begin{enumerate}
\item \emph{Typing.}
\[
\stype{\Gamma[\lvar{x} \mapsto \beta]}{\spec{\lvar{x}}}{\boolS}
\]

\item \emph{Witness existence (the cast image is a model).}
For every valuation $\Delta : \sop{\mathcal{V}} \rightharpoonup \SDom$ and every set $X$ such that
\[
\interpS{\aterm{x}{\Delta}} = \denval{X}[\alpha],
\]
there exists a set $X^!$ such that
\[
\appZ \zop{\mathbb{c}}(X) = X^!
\qquad\text{and}\qquad
\interpS{\aterm{\spec{\lvar{x}}}{\Delta[\lvar{x} \mapsto \denval{X^!}[\beta]]}}
= \denval{\topZ}[\boolS].
\]

\item \emph{Definedness and soundness on typed targets.}
For every valuation $\Delta$, every set $X$ with $\interpS{\aterm{x}{\Delta}} = \denval{X}[\alpha]$, and every typed value $\denval{Y}[\beta]$ (notice that the type carried by $Y$ is $\beta$) the denotation
\(
\interpS{\aterm{\spec{\lvar{x}}}{\Delta[\lvar{x} \mapsto \denval{Y}[\beta]]}}
\)
is defined; and the following additionally holds:
\[
\interpS{\aterm{\spec{\lvar{x}}}{\Delta[\lvar{x} \mapsto \denval{Y}[\beta]]}} = \denval{\topZ}[\boolS] \;\implies\;
\appZ \zop{\mathbb{c}}(X) = Y.
\]
\end{enumerate}
\end{theorem}

Item~2 is a completeness statement: under the valuation assigning the cast image $\appZ \zop{\mathbb{c}}(X)$ to $\lvar{x}$, the formula $\spec{\lvar{x}}$ denotes $\topZ$. Item~3 is a soundness statement: among valuations assigning a typed target to $\lvar{x}$, this can happen only for the cast image. Together they show that $\spec{\lvar{x}}$ has a unique typed solution, namely $\appZ \zop{\mathbb{c}}(X)$; and since $\zop{\mathbb{c}}$ is injective, this image in turn determines $X$, which is the precise sense in which $\lvar{x}$ ``is'' the loosened $x$.

\Cref{thm:loosening-correctness} is depicted as a diagram in \Cref{fig:loosening-square}, in which loosening commutes with denotation and the semantic cast.

\begin{remark}
  The actual statement in Lean is slightly stronger than \Cref{thm:loosening-correctness} above: it also threads fresh-name generation, context extension, weakening, and totality of denotation under valuation updates.
  In the proofs that follow, we suppress this administrative layer and keep only the mathematical core.
  Freshness is used only through the standard fact that updating a valuation at a fresh variable does not change the denotation of the original source term.
\end{remark}

Throughout, we do not construct the proof component of denotation triples, often writing $\irrelprf$ instead: building these witnesses---as the Lean development does at every step---would only burden already heavy arguments, and they are routinely recovered from the local context. Their explicit constructions can be found in the Lean proofs. Recall also that primed operators such as $\aop{\eqS}$ denote the abstract-syntax counterparts of the concrete SMT operators (\Cref{sec:smt-syntax}).

\begin{figure}
\begin{center}
\begin{tikzpicture}[
    node/.style={draw=none, circle, anchor=mid},
    cross/.style={midway, draw=none, fill=white, circle, inner sep=0pt, minimum size=8pt},
    node distance=1.5cm and 1.5cm,
    minimum size=2em]
  \begin{scope}
    \node (x)    at (0,0)     {$x$};
    \node (x!)   at (4.8,0)    {$\langle\lvar{x}, \lvar{x}_{\mathrm{spec}} \rangle$};
    \node (X)    at (0,-2.1)   {$\denval{X}[\alpha][\irrelprf]$};
    \node (X!)   at (3.4,-2.1) {$\denval{X\mathrlap{\smash{^{!}}}}[\beta][\irrelprf]$};
    \node (spec) at (6.2,-2.1) {$\denval{\topZ}[\boolS][\irrelprf]$};

    \draw[funarr] (x) to node[midway, above] {$\mathrm{loosen}_{\mathbb{c}}$} (x!);
    \draw[funarr] (x) to node[midway, left] {$\interpS{\aterm{\cdot}{\Delta}}$} (X);
    \draw[relarr, draw=black] (x!) to node[pos=0.45, left] {$\interpS{\aterm{\cdot}{\Delta'}}$} (X!);
    \draw[relarr, draw=black] (x!) to node[pos=0.45, right] {$\interpS{\aterm{\cdot}{\Delta'}}$} (spec);
    \draw[relarr, draw=black] (X) to node[midway, below] {$\appZ \zop{\mathbb{c}}(\cdot)$} (X!);
  \end{scope}
\end{tikzpicture}
\end{center}
\caption{The loosening diagram. Solid arrows depict assumptions, dashed arrows denote conclusions, and $\Delta' \coloneq \Delta[\lvar{x}\mapsto\denval{X^!}[\beta]]$.}
\label{fig:loosening-square}
\end{figure}

\subsection{Case \texttt{refl}}

\begin{lemma}[Correctness of \texttt{refl}][loosenAux\_prf\_exact.refl][SMT/Reasoning/Basic/LoosenAuxExact/Refl.html\#loosenAux_prf_exact.refl]
Assume $\stype{\Gamma}{x}{\alpha}$ and let $\mathbb{c} : \castpath{\alpha}{\alpha}$. Then the three conclusions of \Cref{thm:loosening-correctness} hold with $\beta = \alpha$ and semantic cast $\zop{\mathbb{c}} = \idZ{\interpZ{\alpha}}$ corresponding to \eqref{eq:castZF-refl}.
\end{lemma}

\begin{proof}
We verify the three conclusions of \Cref{thm:loosening-correctness}: typing of the specification; existence of a cast-image assignment under which it denotes $\topZ$; and definedness on every typed assignment, with uniqueness of the satisfying one.
By construction, $\spec{\lvar{x}} \;\triangleq\; \lvar{x} \eqS x$ (see \Cref{case:refl}).
Fix a valuation $\Delta$ and a set $X$ with $\interpS{\aterm{x}{\Delta}} = \denval{X}$; semantic typing (\Cref{thm:b-denote-type-det}) gives $\tau_X = \alpha$. As $\lvar{x}$ is fresh, updating $\Delta$ at $\lvar{x}$ leaves the denotation of $x$ unchanged:
\begin{equation}\label{eq:refl_fresh}
  \interpS{\aterm{x}{\Delta[\lvar{x} \mapsto \denval{Y}[\alpha]]}} = \denval{X}[\alpha]
  \qquad\text{for any } Y. \tag{\ensuremath{\star}}
\end{equation}
Hence, for any typed $Y$ carrying the SMT type $\alpha$, writing $\Delta' \coloneq \Delta[\lvar{x} \mapsto \denval{Y}[\alpha]]$,
\begin{align*}
  \interpS{\aterm{\spec{\lvar{x}}}{\Delta'}}
  &\;\;=\;\; \interpS{\aterm{\lvar{x} \eqS x}{\Delta'}} & \text{by construction} \\
  &\;\;=\;\; \interpS{\aterm{\lvar{x}}{\Delta'} \aop{\eqS} \aterm{x}{\Delta'}} & \text{by definition of \(\interpS{\cdot}\)}\\
  &\;\;=\;\; \denval{Y \eqZ X}[\boolS] & \text{by \eqref{eq:refl_fresh}}
\end{align*}
where the set-level equality $Y \eqZ X$ denotes $\topZ$ exactly when $Y = X$ and $\botZ$ otherwise. The cast image is $X^! \coloneq \idZ{\interpZ{\alpha}}(X) = X$.

Taking $Y = X^! = X$ therefore reduces this equality to \(\topZ\), which is item~2, while the displayed equivalence is precisely the definedness and soundness of item~3. Here the loosening diagram (\Cref{fig:loosening-square}) degenerates, its semantic cast collapsing to the identity, as shown below.

\begin{center}
\begin{tikzpicture}[
    node/.style={draw=none, circle, anchor=mid},
    node distance=1.5cm and 1.5cm,
    minimum size=2em]
  \node (x)    at (0,0)     {$x$};
  \node (x!)   at (4.8,0)    {$\langle\lvar{x}, \spec{\lvar{x}} \rangle$};
  \node (X)    at (0,-2.1)   {$\denval{X}[\alpha][\irrelprf]$};
  \node (X!)   at (3.6,-2.1) {$\denval{X}[\alpha][\irrelprf]$};
  \node (spec) at (6,-2.1) {$\denval{\topZ}[\boolS][\irrelprf]$};

  \draw[funarr] (x) to node[midway, above] {$\mathrm{loosen}_{\mathbb{c}}$} (x!);
  \draw[funarr] (x) to node[midway, left] {$\interpS{\aterm{\cdot}{\Delta}}$} (X);
  \draw[relarr, draw=black] (x!) to node[pos=0.45, left] {$\interpS{\aterm{\cdot}{\Delta'}}$} (X!);
  \draw[relarr, draw=black] (x!) to node[pos=0.45, right] {$\interpS{\aterm{\cdot}{\Delta'}}$} (spec);
  \draw[relarr, draw=black] (X) to node[midway, below] {$\idZ{\interpZ{\alpha}}$} (X!);
\end{tikzpicture}
\end{center}
\end{proof}

\subsection{Case \texttt{opt}}

\begin{lemma}[Correctness of \texttt{opt}][loosenAux\_prf\_exact.opt][SMT/Reasoning/Basic/LoosenAuxExact/Opt.html\#loosenAux_prf_exact.opt]
\label{lem:loosening-opt}
Assume $\stype{\Gamma}{x}{\optS\,\alpha}$, let $\mathbb{c}_\alpha : \castpath{\alpha}{\alpha'}$, and hence obtain $\mathbb{c} : \castpath{\optS\,\alpha}{\optS\,\alpha'}$. Then \Cref{thm:loosening-correctness} holds with target type $\optS\,\alpha'$ and semantic cast \eqref{eq:castZF-opt}.
\end{lemma}

\begin{proof}
We verify the three conclusions of \Cref{thm:loosening-correctness}: typing of the specification; existence of a cast-image assignment under which it denotes $\topZ$; and definedness on every typed assignment, with uniqueness of the satisfying one.
Fix $\Delta$ with $\interpS{\aterm{x}{\Delta}} = \denval{X}$; then by \Cref{thm:b-denote-type-det}, $\tau_X = \optS\,\alpha$, so $X$ is either $\noneZ$ or $\someZ\ Y$ with $Y \in \interpZ{\alpha}$.
The semantic option cast \eqref{eq:castZF-opt} acts as $\zop{\mathbb{c}}(\noneZ) = \noneZ$ and $\zop{\mathbb{c}}(\someZ\ Y) = \someZ\ \appZ\zop{\mathbb{c}_\alpha}(Y)$, loosening $Y$ by the inner cast. The definition branches on the syntactic shape of the term $x$:

\begin{description}
  \item[If $x = \noneS_{\alpha}\,$:] take $X^! \coloneq \noneZ$.
  \item[If $x = \someS\ y\,$:] the specification is
    \[
      \spec{\lvar{x}} \triangleq \existsS \lvar{y}:\alpha'\cdot\ (\lvar{x} \eqS \someS\ \lvar{y}) \andS \spec{\lvar{y}},
    \]
    where $\langle \lvar{y}, \spec{\lvar{y}} \rangle \coloneq \mathsf{loosen}_{\mathbb{c}_\alpha}(y)$
    is the recursive loosening of $y$ (\Cref{case:opt}). Writing $\interpS{\aterm{y}{\Delta}} = \denval{Y}[\alpha]$, so that $X = \someZ\ Y$, the induction hypothesis on $y$ gives $Y^!$ with $\appZ\zop{\mathbb{c}_\alpha}(Y) = Y^!$ and
    \[
      \interpS{\aterm{\spec{\lvar{y}}}{\Delta'}} = \denval{\topZ}[\boolS]
    \]
    where $\Delta' \coloneq \Delta[\lvar{y} \mapsto \denval{Y^!}[\alpha']]$.
    Set $X^! \coloneq \someZ\ Y^!$. By \eqref{eq:castZF-opt},
    \[
      \appZ\zop{\mathbb{c}}(X) = \someZ\ \appZ\zop{\mathbb{c}_\alpha}(Y) = X^!.
    \]
    Write $\Delta_{\lvar{x}} \coloneq \Delta[\lvar{x} \mapsto \denval{X^!}[\optS\,\alpha']]$.
    \begin{fact}
    Under the valuation $\Delta_{\lvar{x}}$, which assigns the cast image to $\lvar{x}$, the loosened specification $\spec{\lvar{x}}$ denotes $\topZ$:
    \[
      \interpS{\aterm{\spec{\lvar{x}}}{\Delta_{\lvar{x}}}} = \denval{\topZ}[\boolS].
    \]
    \end{fact}
    \begin{subproof}
    Let $\psi \coloneq (\lvar{x} \eqS \someS\ \lvar{y}) \andS \spec{\lvar{y}}$ be the existential body. Take $Y^!$ as the witness and set $\Delta'' \coloneq \Delta_{\lvar{x}}[\lvar{y} \mapsto \denval{Y^!}[\alpha']]$. Under $\Delta''$, both $\lvar{x}$ and $\someS\ \lvar{y}$ denote $\someZ\ Y^!$, so the equality conjunct denotes $\topZ$. The recursive conjunct also denotes $\topZ$ by the induction hypothesis, since $\Delta''$ agrees with $\Delta'$ on the free variables of $\spec{\lvar{y}}$. Thus, $\interpS{\aterm{\psi}{\Delta''}} \;=\; \topZ \andZ \topZ \;=\; \topZ$, and the existential specification denotes $\topZ$.
    \end{subproof}
  \item[Otherwise:] the definition emits the dynamic test $\iteS (x \eqS \noneS_{\alpha})\,(\cdots)\,(\cdots)$. If $X = \noneZ$, it reduces to the first case. Otherwise $X\neq\noneZ$. Since $X\in\optZ(\interpZ{\alpha})$, there is a unique $Y_0\in\interpZ{\alpha}$ such that $X=\someZ\,Y_0$. Consequently,
  \[
    \interpS{\aterm{\theS x}{\Delta''}} = \denval{Y_0}[\alpha],
  \]
  and the argument reduces to the syntactic $\someS$ case with $X^! = \someZ\,Y_0^!$.
\end{description}
In every branch $\appZ\zop{\mathbb{c}}(X) = X^!$, which is item~2 of \Cref{thm:loosening-correctness}.

For item~3, a typed target $\denval{W}[\optS\,\alpha']$ is either $\noneZ$ or $\someZ\ U$; in both cases the equalities and the recursive $\spec{\lvar{y}}$ are defined, hence so is $\spec{\lvar{x}}$, and if it holds then the same case analysis identifies $W$ with $\appZ\zop{\mathbb{c}}(X)$.

\begin{center}
\begin{tikzpicture}[
    node/.style={draw=none, circle, anchor=mid},
    node distance=1.5cm and 1.5cm,
    minimum size=2em]
  \node (x)    at (0,0)      {$x$};
  \node (x!)   at (4.7,0)    {$\langle\lvar{x}, \spec{\lvar{x}} \rangle$};
  \node (X)    at (0,-2.3)   {$\left\langle\begin{smallmatrix}\someZ\ Y\\[2pt] \optS\,\alpha\end{smallmatrix}\right\rangle$};
  \node (X!)   at (3.4,-2.3) {$\left\langle\begin{smallmatrix}\someZ\ Y^!\\[2pt] \optS\,\alpha'\end{smallmatrix}\right\rangle$};
  \node (spec) at (6,-2.3)   {$\left\langle\begin{smallmatrix}\topZ\\[2pt] \boolS\end{smallmatrix}\right\rangle$};

  \draw[funarr] (x) to node[midway, above] {$\mathrm{loosen}_{\mathbb{c}}$} (x!);
  \draw[funarr] (x) to node[midway, left] {$\interpS{\aterm{\cdot}{\Delta}}$} (X);
  \draw[relarr, draw=black] (x!) to node[pos=0.45, left] {$\interpS{\aterm{\cdot}{\Delta_{\lvar{x}}}}$} (X!);
  \draw[relarr, draw=black] (x!) to node[pos=0.45, right] {$\interpS{\aterm{\cdot}{\Delta_{\lvar{x}}}}$} (spec);
  \draw[relarr, draw=black] (X) to node[midway, below] {$\zop{\mathbb{c}}$} (X!);
\end{tikzpicture}
\end{center}

Here the loosening diagram (\Cref{fig:loosening-square}) specializes to the $\someS$ constructor: its semantic cast carries $\someZ\ Y$ to $\someZ\ Y^!$, as shown above.
\end{proof}

\subsection{Case \texttt{pair}}

\begin{lemma}[Correctness of \texttt{pair}][loosenAux\_prf\_exact.pair][SMT/Reasoning/Basic/LoosenAuxExact/Pair.html\#loosenAux_prf_exact.pair]
\label{lem:loosening-pair}
Assume $\stype{\Gamma}{x}{\alpha \timesS \beta}$, let $\mathbb{c}_\alpha : \castpath{\alpha}{\alpha'}$, $\mathbb{c}_\beta : \castpath{\beta}{\beta'}$, and let $\mathbb{c} \coloneq (\mathbb{c}_\alpha, \mathbb{c}_\beta) : \castpath{(\alpha \timesS \beta)}{(\alpha' \timesS \beta')}$. \Cref{thm:loosening-correctness} holds with target type $\alpha' \timesS \beta'$ and semantic cast \eqref{eq:castZF-pair}, the product of the component casts.
\end{lemma}

\begin{proof}
We verify the three conclusions of \Cref{thm:loosening-correctness}: typing of the specification; existence of a cast-image assignment under which it denotes $\topZ$; and definedness on every typed assignment, with uniqueness of the satisfying one.
Fix a valuation $\Delta$ and a set $X$ with $\interpS{\aterm{x}{\Delta}} = \denval{X}$; semantic typing (\Cref{thm:b-denote-type-det}) gives $\tau_X = \alpha \timesS \beta$, so $X = (a, b)$ is a Kuratowski pair with $a \in \interpZ{\alpha}$ and $b \in \interpZ{\beta}$, and the two projections denote the components,
\[
  \interpS{\aterm{\fstS x}{\Delta}} = \denval{a}[\alpha]
  \qquad\text{and}\qquad
  \interpS{\aterm{\sndS x}{\Delta}} = \denval{b}[\beta].
\]
Applying the induction hypothesis to $\fstS x$ and $\sndS x$ yields loosened components $a^!, b^!$ with
\[
  \appZ\zop{\mathbb{c}_\alpha}(a) = a^!, \qquad
  \appZ\zop{\mathbb{c}_\beta}(b) = b^!,
\]
each recursive specification holding at its cast image: writing $\Delta_{\lvar{fst}} \coloneq \Delta[\lvar{fst} \mapsto \denval{a^!}[\alpha']]$ and $\Delta_{\lvar{snd}} \coloneq \Delta[\lvar{snd} \mapsto \denval{b^!}[\beta']]$,
\[
  \interpS{\aterm{\spec{\lvar{fst}}}{\Delta_{\lvar{fst}}}} = \denval{\topZ}[\boolS]
  \qquad\text{and}\qquad
  \interpS{\aterm{\spec{\lvar{snd}}}{\Delta_{\lvar{snd}}}} = \denval{\topZ}[\boolS].
\]
Set $X^! \coloneq (a^!, b^!)$. Since the pair cast is the product of the component casts \eqref{eq:castZF-pair},
\begin{equation}\label{eq:pair_castimg}
  \appZ\zop{\mathbb{c}}(X) \stackrel{\eqref{eq:castZF-pair}}{=} \bigl(\appZ\zop{\mathbb{c}_\alpha}(a),\, \appZ\zop{\mathbb{c}_\beta}(b)\bigr) = (a^!, b^!) = X^!.
  \tag{\ensuremath{\star}}
\end{equation}

\begin{fact}
Under the valuation $\Delta_{\lvar{x}} \coloneq \Delta[\lvar{x} \mapsto \denval{X^!}[\alpha' \timesS \beta']]$, which assigns the cast image to $\lvar{x}$, the loosened pair specification $\spec{\lvar{x}}$ denotes $\topZ$:
\[
  \interpS{\aterm{\spec{\lvar{x}}}{\Delta_{\lvar{x}}}} = \denval{\topZ}[\boolS].
\]
\end{fact}
\begin{subproof}
The pair specification (\Cref{case:pair}) is
\[
  \spec{\lvar{x}} \triangleq \existsS \lvar{fst}:\alpha',\, \lvar{snd}:\beta'\cdot\ (\lvar{x} \eqS \pairS{ \lvar{fst}}{\lvar{snd}}) \andS \spec{\lvar{fst}} \andS \spec{\lvar{snd}}.
\]
As in the \texttt{opt} case, it suffices to exhibit two witnesses under which the body is denoted by $\topZ$. Take the pair $(a^!, b^!)$ and $\Delta'' \coloneq \Delta_{\lvar{x}}[\lvar{fst} \mapsto \denval{a^!}[\alpha'],\, \lvar{snd} \mapsto \denval{b^!}[\beta']]$.
Then, since $\lvar{x}$ and $\pairS{ \lvar{fst}}{\lvar{snd}}$ are both denoted by $(a^!, b^!) = X^!$ under $\Delta''$, and $\Delta''$ agrees with $\Delta_{\lvar{fst}}$ and $\Delta_{\lvar{snd}}$ on the free variables of $\spec{\lvar{fst}}$ and $\spec{\lvar{snd}}$ respectively (neither contains $\lvar{x}$),
\begin{fitalign}
  \interpS{\aterm{\lvar{x} \eqS \pairS{ \lvar{fst}}{\lvar{snd}}}{\Delta''}} &\;=\; \denval{(a^!, b^!) \eqZ (a^!, b^!)}[\boolS][\irrelprf] = \denval{\topZ}[\boolS][\irrelprf],\\
  \interpS{\aterm{\spec{\lvar{fst}}}{\Delta''}} &\;=\; \denval{\topZ}[\boolS][\irrelprf],\\
  \interpS{\aterm{\spec{\lvar{snd}}}{\Delta''}} &\;=\; \denval{\topZ}[\boolS][\irrelprf],
\end{fitalign}
Thus, $\interpS{\aterm{\psi}{\Delta''}} = \topZ \andZ \topZ \andZ \topZ = \topZ$, and therefore $\interpS{\aterm{\spec{\lvar{x}}}{\Delta_{\lvar{x}}}} = \notZ\botZ = \topZ$.
\end{subproof}

Together with \eqref{eq:pair_castimg} this is item~2. For item~3, a typed target $\denval{Y}[\alpha' \timesS \beta']$ decomposes as $Y = (u, v)$; the pair equality is boolean and the recursive specifications $\spec{\lvar{fst}}, \spec{\lvar{snd}}$ are defined on typed data, so $\spec{\lvar{x}}$ is defined. If it holds at $Y$, recursive soundness identifies the witnesses with the components, $u = \appZ\zop{\mathbb{c}_\alpha}(a)$ and $v = \appZ\zop{\mathbb{c}_\beta}(b)$, whence $Y = (u, v) = (a^!, b^!) = X^! = \appZ\zop{\mathbb{c}}(X)$ by \eqref{eq:pair_castimg}.

Here the loosening diagram (\Cref{fig:loosening-square}) specializes to the product, its semantic cast $\zop{\mathbb{c}} = (\zop{\mathbb{c}_\alpha}, \zop{\mathbb{c}_\beta})$ acting componentwise:

\begin{center}
\begin{tikzpicture}[node distance=1.5cm and 1.5cm, minimum size=2em]
  \node (x)    at (0,0)      {$x$};
  \node (x!)   at (5.0,0)    {$\langle\lvar{x}, \spec{\lvar{x}} \rangle$};
  \node (X)    at (0,-2.4)   {$\left\langle\begin{smallmatrix}(a,\, b)\\[2pt] \alpha \timesS \beta\end{smallmatrix}\right\rangle$};
  \node (X!)   at (3.7,-2.4) {$\left\langle\begin{smallmatrix}(a^!,\, b^!)\\[2pt] \alpha' \timesS \beta'\end{smallmatrix}\right\rangle$};
  \node (spec) at (6.7,-2.4) {$\left\langle\begin{smallmatrix}\topZ\\[2pt] \boolS\end{smallmatrix}\right\rangle$};
  \draw[funarr] (x) to node[midway, above] {$\mathrm{loosen}_{(\mathbb{c}_\alpha, \mathbb{c}_\beta)}$} (x!);
  \draw[funarr] (x) to node[midway, left] {$\interpS{\aterm{\cdot}{\Delta}}$} (X);
  \draw[relarr, draw=black] (x!) to node[pos=0.45, left] {$\interpS{\aterm{\cdot}{\Delta_{\lvar{x}}}}$} (X!);
  \draw[relarr, draw=black] (x!) to node[pos=0.45, right] {$\interpS{\aterm{\cdot}{\Delta_{\lvar{x}}}}$} (spec);
  \draw[relarr, draw=black] (X) to node[midway, below] {$({\zop{\mathbb{c}}}_\alpha, {\zop{\mathbb{c}}}_\beta)$} (X!);
\end{tikzpicture}
\vspace*{-\baselineskip}
\end{center}
\end{proof}

\subsection{Case \texttt{chpred}}

\begin{lemma}[Correctness of \texttt{chpred}][loosenAux\_prf\_exact.chpred][SMT/Reasoning/Basic/LoosenAuxExact/Chpred.html\#loosenAux_prf_exact.chpred]
\label{lem:loosening-chpred}
Assume $\stype{\Gamma}{x}{\alpha \toS \boolS}$, let $\mathbb{c}_\alpha : \castpath{\alpha}{\alpha'}$ and let $\mathbb{c} \coloneq \texttt{chpred}\,\mathbb{c}_\alpha : \castpath{(\alpha \toS \boolS)}{(\alpha' \toS \boolS)}$. \Cref{thm:loosening-correctness} holds with target type $\alpha' \toS \boolS$ and semantic cast $\zop{\mathbb{c}}$ \eqref{eq:castZF-chpred}.
\end{lemma}

\begin{proof}
We verify the three conclusions of \Cref{thm:loosening-correctness}: typing of the specification; existence of a cast-image assignment under which it denotes $\topZ$; and definedness on every typed assignment, with uniqueness of the satisfying one.
Fix a valuation $\Delta$ and a set $X$ with $\interpS{\aterm{x}{\Delta}} = \denval{X}$.
By \Cref{thm:b-denote-type-det}, \(\tau_X = \alpha \toS \boolS\) and \(X \in \mathbb{B}^{\interpZ{\alpha}}\) is a characteristic predicate. The construction (\Cref{case:chpred}) loosens a fresh source variable $\var{z} : \alpha$ by the inner cast, producing $\langle \lvar{z}, \spec{\lvar{z}} \rangle \coloneq \mathsf{loosen}_{\mathbb{c}_\alpha}(\var{z})$, whose specification computes the graph of \(\zop{\mathbb{c}}_\alpha\).

\begin{fact}
For every $Z \in \interpZ{\alpha}$ and $Z^! \in \interpZ{\alpha'}$, the denotation $\Phi_{Z,Z^!}$ of $\spec{\lvar{z}}$ is defined under $\Delta_{\var{z}, \var{z}^!} \coloneq \Delta[\var{z} \mapsto \denval{Z}[\alpha]$, $\lvar{z} \mapsto \denval{Z^!}[\alpha']]$, and
\begin{equation}\label{eq:zspec_chpred}
  \Phi_{Z,Z^!} = \topZ \iff \appZ \zop{\mathbb{c}}_\alpha(Z) = Z^!. \tag{\ensuremath{\star}}
\end{equation}
\end{fact}
\begin{subproof}
The induction is on the cast path, so applying \Cref{thm:loosening-correctness} recursively to the strict subpath $\mathbb{c}_\alpha$ is legitimate. Apply it to the fresh source variable $\var{z} : \alpha$, well typed with $\interpS{\aterm{\var{z}}{\Delta_{\var{z}, \var{z}^!}}} = \denval{Z}[\alpha]$. Item~3 gives that $\spec{\lvar{z}}$ is defined at every typed target $Z^!$ and, by soundness, the forward implication: if $\Phi_{Z,Z^!} = \topZ$ then $Z^! = \appZ \zop{\mathbb{c}}_\alpha(Z)$. Item~2 gives the converse: at the cast image $Z^! = \appZ \zop{\mathbb{c}}_\alpha(Z)$ the specification holds, so $\Phi_{Z,\appZ \zop{\mathbb{c}}_\alpha(Z)} = \topZ$. Together they yield \eqref{eq:zspec_chpred}.
\end{subproof}

Define $X^! \in \mathbb{B}^{\interpZ{\alpha'}}$ as the transport of $X$ along $\zop{\mathbb{c}}_\alpha$:
\[
X^! \coloneq
\begin{cases}
  X \fcompZ {\zop{\mathbb{c}}_\alpha}^{-1} & \text{on } \Ran(\zop{\mathbb{c}}_\alpha),\\
  \botZ & \text{off } \Ran(\zop{\mathbb{c}}_\alpha).
\end{cases}%
\]
This is well-defined since $\zop{\mathbb{c}}_\alpha$ is injective, hence invertible on its range; and $\appZ \zop{\mathbb{c}}(X) = X^!$ by the characteristic-predicate cast \eqref{eq:castZF-chpred}.

The \texttt{chpred} case (\Cref{case:chpred}) defines
\[
  \spec{\lvar{x}} \;\;\triangleq\;\; \lvar{x} \eqS \left(\lamS \lvar{z}:\alpha' \cdot\, \existsS \var{z}:\alpha\cdot\ \appS x(\var{z}) \andS \spec{\lvar{z}}\right),
\] whose $\lamS$-body denotes $X^!$ pointwise.

\begin{fact}
For every $a' \in \interpZ{\alpha'}$, writing $\Delta_{\lvar{z}} \coloneq \Delta[\lvar{z} \mapsto \denval{a'}[\alpha']]$---the source variable $\var{z}$ is bound by the inner existential, hence left unassigned,
\[
  \interpS{\aterm{\existsS \var{z}:\alpha\cdot\ \appS x(\var{z}) \andS \spec{\lvar{z}}}{\Delta_{\lvar{z}}}} = X^!(a').
\]
\end{fact}
\begin{subproof}
If $a' \notin \Ran(\zop{\mathbb{c}}_\alpha)$, then \eqref{eq:zspec_chpred} makes $\spec{\lvar{z}}$ denote $\botZ$ for every source value $a$. Hence the existential denotes $\botZ=X^!(a')$.

If $a' \in \Ran(\zop{\mathbb{c}}_\alpha)$, let $a_0 \in \interpZ{\alpha}$ be the unique preimage satisfying $\appZ \zop{\mathbb{c}}_\alpha(a_0)=a'$; uniqueness follows from injectivity. By \eqref{eq:zspec_chpred}, only $a_0$ can make $\spec{\lvar{z}}$ true. The existential therefore denotes $\topZ$ exactly when $\appZ X(a_0)=\topZ$, which is exactly when $X^!(a')=\topZ$.
\end{subproof}

By the second fact the $\lamS$-body denotes $X^!$, so the (denotation of the) outer equality $\spec{\lvar{x}}$ holds (\ie equals \(\topZ\)) under $\Delta[\lvar{x} \mapsto \denval{X^!}[\alpha' \toS \boolS]]$; with $\appZ \zop{\mathbb{c}}(X) = X^!$ this is item~2.

For item~3, at any typed target $\denval{Y}[\alpha' \toS \boolS]$ the $\lamS$-body has a denotation (it is built from $\appS x(\var{z})$ and $\spec{\lvar{z}}$), and so has $\spec{\lvar{x}}$. It is therefore \(\topZ\) at $Y$ if and only if $Y$ equals the body's denotation $\appZ \zop{\mathbb{c}}(X)$, \ie if and only if $Y = X^!$.

The loosening diagram (\Cref{fig:loosening-square}) here transports a predicate along $\zop{\mathbb{c}}_\alpha$, its semantic cast relabeling the domain by $a' = \appZ \zop{\mathbb{c}}_\alpha(a)$ and defaulting to $\botZ$ off $\Ran(\zop{\mathbb{c}}_\alpha)$:

\begin{center}
\begin{tikzpicture}[node distance=1.5cm and 1.5cm, minimum size=2em]
  \node (x)    at (0,0)      {$x$};
  \node (x!)   at (4.7,0)    {$\langle\lvar{x}, \spec{\lvar{x}} \rangle$};
  \node (X)    at (0,-2.3)   {$\left\langle\begin{smallmatrix}X\\[2pt] \alpha \toS \boolS\end{smallmatrix}\right\rangle$};
  \node (X!)   at (3.5,-2.3) {$\left\langle\begin{smallmatrix}X^!\\[2pt] \alpha' \toS \boolS\end{smallmatrix}\right\rangle$};
  \node (spec) at (6.2,-2.3) {$\left\langle\begin{smallmatrix}\topZ\\[2pt] \boolS\end{smallmatrix}\right\rangle$};
  \draw[funarr] (x) to node[midway, above] {$\mathrm{loosen}_{\mathbb{c}}$} (x!);
  \draw[funarr] (x) to node[midway, left] {$\interpS{\aterm{\cdot}{\Delta}}$} (X);
  \draw[relarr, draw=black] (x!) to node[pos=0.45, left] {$\interpS{\aterm{\cdot}{\Delta_{\lvar{x}}}}$} (X!);
  \draw[relarr, draw=black] (x!) to node[pos=0.45, right] {$\interpS{\aterm{\cdot}{\Delta_{\lvar{x}}}}$} (spec);
  \draw[relarr, draw=black] (X) to node[midway, below] {$\zop{\mathbb{c}}$} (X!);
\end{tikzpicture}
\vspace*{-\baselineskip}
\end{center}
\end{proof}

\subsection{Case \texttt{fun}}

\begin{lemma}[Correctness of \texttt{fun}][loosenAux\_prf\_exact.fun][SMT/Reasoning/Basic/LoosenAuxExact/Fun.html\#loosenAux_prf_exact.fun]
\label{lem:loosening-fun}
Assume $\stype{\Gamma}{x}{\alpha \toS \beta}$ with $\beta \neq \boolS$. Let $\mathbb{c}_\alpha : \castpath{\alpha}{\alpha'}$ and $\mathbb{c}_\beta : \castpath{\beta}{\beta'}$ (thus \(\beta' \neq \boolS\)), and let ${\mathbb{c} \coloneq \texttt{fun}\,\mathbb{c}_\alpha\,\mathbb{c}_\beta : \castpath{(\alpha \toS \beta)}{(\alpha' \toS \beta')}}$.

\Cref{thm:loosening-correctness} holds with target type $\alpha' \toS \beta'$ and semantic cast \eqref{eq:castZF-fun}; moreover, as $\beta \neq \boolS$, item~3 sharpens to a pointwise characterization.
\end{lemma}

\begin{proof}
We verify the three conclusions of \Cref{thm:loosening-correctness}: typing of the specification; existence of a cast-image assignment under which it denotes $\topZ$; and definedness on every typed assignment, with uniqueness of the satisfying one.
Fix a valuation $\Delta$ and a set $X$ with $\interpS{\aterm{x}{\Delta}} = \denval{X}$. So, $\tau_X = \alpha \toS \beta$ and $X \in {\interpZ{\beta}}^{\interpZ{\alpha}}$ is a total function. The component casts $\zop{\mathbb{c}}_\alpha$ and $\zop{\mathbb{c}}_\beta$ are injective by \Cref{lem:semantic-cast-injective}. With $d_{\beta'}$ the canonical default of $\interpZ{\beta'}$ (\Cref{def:smt-default}), the semantic function cast \eqref{eq:castZF-fun} conjugates $X$ by the component casts:

\noindent
\begin{minipage}[t]{0.7\textwidth}
\[
X^! \coloneq
\begin{cases}
\zop{\mathbb{c}}_\beta \fcompZ X \fcompZ {\zop{\mathbb{c}}_\alpha}^{-1} & \text{on } \Ran(\zop{\mathbb{c}}_\alpha),\\
d_{\beta'} & \text{off } \Ran(\zop{\mathbb{c}}_\alpha)
\end{cases}
\qquad\text{and}\qquad \appZ\zop{\mathbb{c}}(X) = X^!,
\]
This definition of $X^!$ is well-defined because $\zop{\mathbb{c}}_\alpha$ is injective and therefore invertible on its range.
\end{minipage}\hfill
\begin{minipage}[t]{0.25\textwidth}
\vspace{0pt}
\centering
\begin{tikzpicture}
  \node (A)  at (0,0)      {$\interpZ{\alpha}$};
  \node (B)  at (2.2,0)    {$\interpZ{\beta}$};
  \node (A') at (0,-2)   {$\interpZ{\alpha'}$};
  \node (B') at (2.2,-2) {$\interpZ{\beta'}$};
  \draw[funarr] (A)  to node[midway,above]{$X$} (B);
  \draw[funarr] (A') to node[midway,below]{$X^!$} (B');
  \draw[funarr] (A)  to node[midway,left]{$\zop{\mathbb{c}}_\alpha$} (A');
  \draw[funarr] (B)  to node[midway,right]{$\zop{\mathbb{c}}_\beta$} (B');
\end{tikzpicture}
\end{minipage}

The specification (\Cref{case:fun}) tests, for each target argument $\lvar{a}$, whether it lies in the image of the domain cast, and accordingly pins the value pointwise or forces the default:
\begin{align*}
  \spec{\lvar{x}} \;\;\triangleq\;\; {} & \forallS \lvar{a} : \alpha'\cdot\, \\
  &\qquad \iteS \left(\existsS \var{a} : \alpha\cdot\, \spec{\lvar{a}}\right) \\
  &\qquad\quad \left(
    \begin{aligned}
      &\forallS \lvar{b} : \beta'\cdot\, ((\appS{\lvar{x}}{(\lvar{a})} \eqS \lvar{b}) \eqS {}\\
      &\qquad (\existsS \var{a} : \alpha,\, \var{b} : \beta\cdot\, (\appS x (\var{a}) \eqS \var{b}) \andS (\spec{\lvar{a}} \andS \spec{\lvar{b}})))
    \end{aligned}
  \right) \\
  &\qquad\quad \left(\mathsf{defaultSpec}_{\beta'}(\appS{\lvar{x}}{(\lvar{a})})\right)
\end{align*}
The default branch relies on the following characterization of $\mathsf{defaultSpec}$.

\begin{fact}
For every SMT type $\gamma$ and term $t$ with $\stype{\Gamma}{t}{\gamma}$, and every \(\Delta\) such that $\interpS{\aterm{t}{\Delta}} = \denval{D}[\gamma]$, the predicate $\mathsf{defaultSpec}_{\gamma}(t)$ is such that
\[
  \interpS{\aterm{\mathsf{defaultSpec}_{\gamma}(t)}{\Delta}} = \denval{\topZ}[\boolS]
  \iff
  D = d_{\gamma},
\]
where $d_{\gamma}$ is defined by $\interpS{\mathsf{default}(\gamma)} = \denval{d_{\gamma}}[\gamma]$, so $d_{\gamma} \in \interpZ{\gamma}$ is the canonical default value.
\end{fact}
\begin{subproof}
Proceed by structural induction on $\gamma$, following the definition of $\mathsf{defaultSpec}$ (\Cref{case:fun}). The cases $\intS$, $\boolS$, and $\unitS$ are immediate. For $\optS\,\xi$, the defining equality holds exactly when $D=\noneZ=d_{\optS\,\xi}$. For $\xi\timesS\kappa$, the two induction hypotheses identify the components of $D$ with $d_\xi$ and $d_\kappa$, hence $D=(d_\xi,d_\kappa)=d_{\xi\timesS\kappa}$. Finally, for $\xi\toS\kappa$, the induction hypothesis applies pointwise, so the universal specification denotes $\topZ$ exactly when $D$ is the constant map with value $d_\kappa$, namely $d_{\xi\toS\kappa}$.
\end{subproof}

We moreover show that, for a typed target $\denval{Y}[\alpha' \toS \beta']$, the specification holds at $Y$ precisely when $Y = X^!$ pointwise.

\begin{fact}
For every $a' \in \interpZ{\alpha'}$, the $\iteS$-body of $\spec{\lvar{x}}$ at $\lvar{a} \mapsto \denval{a'}[\alpha']$ is denoted by $\topZ$ if and only if $\appZ Y(a') = X^!(a')$.
\end{fact}
\begin{subproof}
If $a' \notin \Ran(\zop{\mathbb{c}}_\alpha)$, then by recursive soundness (item~3 for $\spec{\lvar{a}}$) no source argument $\var{a}$ satisfies $\spec{\lvar{a}}$ at $a'$, so the guard $\existsS \var{a}\cdot\ \spec{\lvar{a}}$ is $\botZ$ and the $\iteS$ takes its default branch $\mathsf{defaultSpec}_{\beta'}(\appS{\lvar{x}}{(\lvar{a})})$. By the previous fact, it is denoted by $\topZ$ exactly when $\appZ Y(a') = d_{\beta'} = X^!(a')$.

If $a' = \appZ \zop{\mathbb{c}}_\alpha(a)$ for the preimage $a$---unique by injectivity of $\zop{\mathbb{c}}_\alpha$---recursive completeness makes the guard $\topZ$, so the $\iteS$ takes its positive branch, which states
\[
  \appZ Y(a') = b' \iff \interpS{\aterm{\existsS \var{a},\var{b}\cdot\ (\appS x(\var{a}) \eqS \var{b}) \andS \spec{\lvar{a}} \andS \spec{\lvar{b}}}{\Delta}} = \denval{\topZ}[\boolS].
\]
By recursive soundness, the only admissible source argument is $a$, and since $X$ is a function, the only source result is $\appZ X(a)$. The recursive characterization of $\spec{\lvar{b}}$ then forces
\[
  b' = \appZ \zop{\mathbb{c}}_\beta(\appZ X(a)) = X^!(a').
\]
So, the denotation of the positive branch of the $\iteS$-specification holds exactly when $\appZ Y(a') = X^!(a')$.
\end{subproof}

Since $\spec{\lvar{x}}$ is the universal quantification $\forallS \lvar{a}\cdot(\cdots)$ over these bodies, it is defined throughout and, by functional extensionality, holds at $Y$ if and only if $\appZ Y(a') = X^!(a')$ for every $a'$, \ie $Y = X^!$. Taking $Y = X^! = \appZ\zop{\mathbb{c}}(X)$ gives item~2, and the biconditional is item~3 in its sharpened, pointwise form---available here because $\beta \neq \boolS$, so the target is a genuine function space.

The loosening diagram (\Cref{fig:loosening-square}) specializes with the function conjugation as its semantic cast:

\begin{center}
\begin{tikzpicture}[node distance=1.5cm and 1.5cm, minimum size=2em]
  \node (x)    at (0,0)      {$x$};
  \node (x!)   at (4.7,0)    {$\langle\lvar{x}, \spec{\lvar{x}} \rangle$};
  \node (X)    at (0,-2.3)   {$\left\langle\begin{smallmatrix}X\\[2pt] \alpha \toS \beta\end{smallmatrix}\right\rangle$};
  \node (X!)   at (3.5,-2.3) {$\left\langle\begin{smallmatrix}X^!\\[2pt] \alpha' \toS \beta'\end{smallmatrix}\right\rangle$};
  \node (spec) at (6.2,-2.3) {$\left\langle\begin{smallmatrix}\topZ\\[2pt] \boolS\end{smallmatrix}\right\rangle$};
  \draw[funarr] (x) to node[midway, above] {$\mathrm{loosen}_{\mathbb{c}}$} (x!);
  \draw[funarr] (x) to node[midway, left] {$\interpS{\aterm{\cdot}{\Delta}}$} (X);
  \draw[relarr, draw=black] (x!) to node[pos=0.45, left] {$\interpS{\aterm{\cdot}{\Delta_{\lvar{x}}}}$} (X!);
  \draw[relarr, draw=black] (x!) to node[pos=0.45, right] {$\interpS{\aterm{\cdot}{\Delta_{\lvar{x}}}}$} (spec);
  \draw[relarr, draw=black] (X) to node[midway, below] {$\zop{\mathbb{c}}$} (X!);
\end{tikzpicture}
\end{center}
\end{proof}

\subsection{Case \texttt{graph}}

\begin{lemma}[Correctness of \texttt{graph}][loosenAux\_prf\_exact.graph][SMT/Reasoning/Basic/LoosenAuxExact/Graph.html\#loosenAux_prf_exact.graph]
\label{lem:loosening-graph}
Assume $\stype{\Gamma}{x}{\alpha \toS \optS\,\beta}$. Let $\mathbb{c}_\alpha : \castpath{\alpha}{\alpha'}$ and $\mathbb{c}_\beta : \castpath{\beta}{\beta'}$, and let $\mathbb{c} \coloneq \texttt{graph}\,\mathbb{c}_\alpha\,\mathbb{c}_\beta : \castpath{(\alpha \toS \optS\,\beta)}{((\alpha' \timesS \beta') \toS \boolS)}$.

\Cref{thm:loosening-correctness} holds with target type $(\alpha' \timesS \beta') \toS \boolS$ and semantic cast \(\bigl(\graphcastZ{\zop{\mathbb{c}_\alpha}}{\zop{\mathbb{c}_\beta}}\bigr)\).
\end{lemma}

\begin{proof}
We verify the three conclusions of \Cref{thm:loosening-correctness}: typing of the specification; existence of a cast-image assignment under which it denotes $\topZ$; and definedness on every typed assignment, with uniqueness of the satisfying one.
Fix a valuation $\Delta$ and a set $X$ with $\interpS{\aterm{x}{\Delta}} = \denval{X}$. By typing, \(\tau_X = \alpha \toS \optS\,\beta\) and $X \in (\optZ\,\interpZ{\beta})^{\interpZ{\alpha}}$ is a partial function. The \texttt{graph} case (\Cref{case:graph}) loosens a fresh source pair $\var{z} : \alpha \timesS \beta$ through the pair cast $\mathbb{c}_{\alpha \timesS \beta}$, yielding $\langle \lvar{z}, \spec{\lvar{z}} \rangle$ whose specification characterizes the pair cast.

\begin{fact}
For $u=(a,b) \in \interpZ{\alpha \timesS \beta}$ and $v=(a',b') \in \interpZ{\alpha' \timesS \beta'}$, the denotation $\Phi_{u,v}$ of $\spec{\lvar{z}}$ under $\Delta[\var{z} \mapsto \denval{u}[\alpha \timesS \beta]$, $\lvar{z} \mapsto \denval{v}[\alpha' \timesS \beta']]$ is such that
\begin{equation}\label{eq:zspec_graph}
  \Phi_{u,v} = \topZ
  \iff
  \appZ \zop{\mathbb{c}_\alpha}(a)=a' \;\land\; \appZ \zop{\mathbb{c}_\beta}(b)=b'.
  \tag{\ensuremath{\star}}
\end{equation}
\end{fact}
\begin{subproof}
The recursive call is $\mathsf{loosen}_{\mathbb{c}_{\alpha \timesS \beta}}(\var{z})$, so \Cref{lem:loosening-pair} applies. Its item~3 gives definedness of the denotation of $\spec{\lvar{z}}$ and, by soundness, the forward implication: a typed $v$ with $\Phi_{u,v} = \topZ$ equals the pair-cast image of $u$, whose components are $\appZ \zop{\mathbb{c}_\alpha}(a)$ and $\appZ \zop{\mathbb{c}_\beta}(b)$. Its item~2 gives the converse, that the denotation of $\spec{\lvar{z}}$ is \(\topZ\) at that cast image. Together they yield \eqref{eq:zspec_graph}.
\end{subproof}

Define the target relation as the following characteristic predicate on $\interpZ{\alpha' \timesS \beta'}$:
\[
X^! \coloneq \zop{\lambda} (a',b') \mapsto
\begin{cases}
\topZ & \text{if } \eqref{eq:x!_graph_def},\\
\botZ & \text{otherwise,}
\end{cases}
\]
where \eqref{eq:x!_graph_def} is the condition
\begin{equation}\label{eq:x!_graph_def}
  a' \in \Ran(\zop{\mathbb{c}_\alpha}),\quad b' \in \Ran(\zop{\mathbb{c}_\beta}),\quad\text{and}\quad
  \appZ X\big(\appZ {\zop{\mathbb{c}_\alpha}}^{-1}(a')\big)=\someZ\ \appZ {\zop{\mathbb{c}_\beta}}^{-1}(b'). \tag{\ensuremath{\star\star}}
\end{equation}
Injectivity of the component casts (\Cref{lem:semantic-cast-injective}) makes the inverses well-defined on their ranges, so $X^! \in \interpZ{(\alpha' \timesS \beta') \toS \boolS}$; this is precisely $\appZ \graphcastZ{\zop{\mathbb{c}_\alpha}}{\zop{\mathbb{c}_\beta}}(X)$, the graph cast \eqref{eq:castZF-graph} applied to $X$.

The \texttt{graph} case defines
\[
  \spec{\lvar{x}} \triangleq \lvar{x} \eqS \lamS \lvar{z}:\alpha' \timesS \beta' \cdot\, \existsS \var{z}:\alpha \timesS \beta\cdot\ (\appS x(\fstS\var{z}) \eqS \someS(\sndS\var{z})) \andS \spec{\lvar{z}},
\] whose $\lamS$-body denotes $X^!$ pointwise.

\begin{fact}
For every $v=(a',b') \in \interpZ{\alpha' \timesS \beta'}$, the existential body is denoted by $X^!(v)$, \ie
\[
  \interpS{\aterm{\existsS \var{z}:\alpha \timesS \beta\cdot\ (\appS x(\fstS\var{z}) \eqS \someS(\sndS\var{z})) \andS \spec{\lvar{z}}}{\Delta[\lvar{z} \mapsto \denval{v}]}} = X^!(v).
\]
\end{fact}
\begin{subproof}
If $a' \notin \Ran(\zop{\mathbb{c}_\alpha})$ or $b' \notin \Ran(\zop{\mathbb{c}_\beta})$, the right-hand side of \eqref{eq:zspec_graph} fails for every $u$, so $\Phi_{u,v} \neq \topZ$ throughout; the existential is $\botZ$, matching $X^!(v) = \botZ$.

Otherwise, let $a = {\zop{\mathbb{c}_\alpha}}^{-1}(a')$ and $b = {\zop{\mathbb{c}_\beta}}^{-1}(b')$ be the unique pre-images. The (denotation of the) existential body at $v$ is $\topZ$ if and only if $\appZ X(a)=\someZ\ b$. Indeed, if $\appZ X(a)=\someZ\ b$ then $u=(a,b)$ is a witness: its first conjunct holds because $\fstZ u = a$, $\sndZ u = b$, and $\appZ X(a)=\someZ\ b$, and (the denotation of) $\spec{\lvar{z}}$ holds at $(u,v)$ by \eqref{eq:zspec_graph}.

Conversely, any witness $u=(a_0,b_0)$ gives $\appZ X(a_0)=\someZ\ b_0$ from the first conjunct and $\zop{\mathbb{c}_\alpha}(a_0)=a',\ \zop{\mathbb{c}_\beta}(b_0)=b'$ from \eqref{eq:zspec_graph}, hence $a_0=a$ and $b_0=b$ by injectivity.

Thus, the body is $\topZ$ exactly on the condition \eqref{eq:x!_graph_def}, \ie it denotes $X^!(v)$.
\end{subproof}

Consequently, under $\Delta_{\lvar{x}} \coloneq \Delta[\lvar{x} \mapsto \denval{X^!}[(\alpha' \timesS \beta') \toS \boolS][\irrelprf]]$, the outer equality $\spec{\lvar{x}}$ holds (\ie is denoted by \(\topZ\)), which is item~2.

For item~3, at any typed target $\denval{Y}[(\alpha' \timesS \beta') \toS \boolS][\irrelprf]$, the $\lamS$-body has a denotation, and so has $\spec{\lvar{x}}$.
If the equality $\lvar{x} \eqS \lamS(\cdots)$ holds at $Y$, then $Y$ equals the $\lamS$-body's denotation $X^! = \appZ \graphcastZ{\zop{\mathbb{c}_\alpha}}{\zop{\mathbb{c}_\beta}}(X)$, that is $(X,Y)$ lies in the graph cast.

This is the one shape-changing rule: the loosening diagram (\Cref{fig:loosening-square}) turns the functional encoding of a partial function into the relational encoding of its graph.

\begin{center}
\begin{tikzpicture}[node distance=1.5cm and 1.5cm, minimum size=2em]
  \node (x)    at (0,0)      {$x$};
  \node (x!)   at (5.7,0)    {$\langle\lvar{x}, \spec{\lvar{x}} \rangle$};
  \node (X)    at (0,-2.5)   {$\left\langle\begin{smallmatrix}X\\[2pt] \alpha \toS \optS\,\beta\end{smallmatrix}\right\rangle$};
  \node (X!)   at (4.3,-2.5) {$\left\langle\begin{smallmatrix}X^!\\[2pt] (\alpha' \timesS \beta') \toS \boolS\end{smallmatrix}\right\rangle$};
  \node (spec) at (7.7,-2.5) {$\left\langle\begin{smallmatrix}\topZ\\[2pt] \boolS\end{smallmatrix}\right\rangle$};
  \draw[funarr] (x) to node[midway, above] {$\mathrm{loosen}_{\graphcastZ{\zop{\mathbb{c}_\alpha}}{\zop{\mathbb{c}_\beta}}}$} (x!);
  \draw[funarr] (x) to node[midway, left] {$\interpS{\aterm{\cdot}{\Delta}}$} (X);
  \draw[relarr, draw=black] (x!) to node[pos=0.45, left] {$\interpS{\aterm{\cdot}{\Delta_{\lvar{x}}}}$} (X!);
  \draw[relarr, draw=black] (x!) to node[pos=0.45, right] {$\interpS{\aterm{\cdot}{\Delta_{\lvar{x}}}}$} (spec);
  \draw[relarr, draw=black] (X) to node[midway, below] {$\graphcastZ{\zop{\mathbb{c}_\alpha}}{\zop{\mathbb{c}_\beta}}$} (X!);
\end{tikzpicture}
\vspace*{-\baselineskip}
\end{center}
\end{proof}

\medskip

Taken together, the six cases exhaust the constructors of a cast path and therefore establish \Cref{thm:loosening-correctness} by induction on $\mathbb{c}$. In each of them the loosening diagram (\Cref{fig:loosening-square}) commutes: the denotation of the loosened term $\lvar{x}$ is the semantic cast $\appZ\zop{\mathbb{c}}(X)$ of the denotation $X$ of the source $x$, and the specification $\spec{\lvar{x}}$ singles out this cast image as its unique typed solution---completeness (item~2) placing the image among the solutions, soundness (item~3) ruling out every other typed target. As the semantic cast is injective (\Cref{lem:semantic-cast-injective}), that image in turn determines $X$, so $\spec{\lvar{x}}$ ties $\lvar{x}$ to $x$ unambiguously: loosening rewrites the \emph{representation} of a value but never the set it denotes. Only the \texttt{graph} case changes the shape of that representation---trading the functional encoding of a partial function for the relational encoding of its graph---while the other five act homomorphically, transporting the value componentwise. This discharges the slogan of \Cref{sec:loosening-rules}: a term may be replaced by its loosened view at a looser type with its denotation provably unchanged.

\section{Conclusion}
\label{sec:loosening-conclusion}

This chapter isolated and resolved the tension identified in \Cref{sec:loosening-motivation}: because a B object is a set while every SMT term carries a fixed type, one and the same set admits several equally faithful SMT representations, and the operators of the encoding must move between them. We called this passage \emph{loosening}---the directional move from a more specific, functional representation to a more general, relational one that leaves the denoted set untouched---and gave a complete order-theoretic and semantic account of it. On types, the \emph{castable} relation $\sqsubseteq$ and the proof-relevant cast-path relation $\castpath{}{}$ were shown to coincide and to induce a decidable partial order in which every type lies in an interval $[\lgf{\alpha}, \mgf{\alpha}]$ between its least (fully functional) and most general (fully relational) forms; the two normal forms constitute a Galois connection $\mathsf{lgf} \dashv \mathsf{mgf}$, and each such interval is a complete lattice. On terms, the $\mathsf{loosen}$ construction realizes this order, returning for every cast path a fresh target variable together with a specification relating it to the source; and \Cref{thm:loosening-correctness} established, rule by rule and against the ZFC semantics, that this specification is sound and definite---it pins the target to the semantic-cast image of the source, and to nothing else.

Type loosening is the first concrete contribution of the encoding developed in this thesis, and the mechanism on which the \beer encoder described in the next chapter (\Cref{ch:beer}) rests: every combinator that must reconcile the types of its operands---the union of two partial functions, the equality of a function with a relation, the membership of a maplet in a set of pairs, or the application of a relation that happens to be a function---does so by loosening one representation into another, secure that the denotation is preserved. Loosening is, by design, one-directional: a function may always be viewed as its graph, but recovering a function from a relation would demand the determinism and totality that loosening deliberately forgets. This asymmetry is exactly what the encoder needs: a single, semantics-preserving coercion toward looser types.
